\documentclass[11pt,a4paper]{ctexart}
\usepackage[margin=2.2cm]{geometry}
\usepackage{amsmath,amssymb,mathtools,bm}
\usepackage{booktabs,longtable,array,tabularx}
\usepackage{enumitem}
\usepackage{microtype}
\usepackage{hyperref}
\usepackage[nameinlink,capitalize]{cleveref}
\usepackage[backend=biber,style=authoryear,maxcitenames=2,maxbibnames=99,doi=true,url=true,isbn=false]{biblatex}
\addbibresource{research_state_control_refs_v4.bib}
\hypersetup{colorlinks=true,linkcolor=blue,citecolor=blue,urlcolor=blue}
\setlist{nosep}
\emergencystretch=3em
\newcommand{\R}{\mathbb{R}}
\newcommand{\E}{\mathbb{E}}
\newcommand{\Pp}{\mathbb{P}}
\newcommand{\Hh}{\mathcal{H}}
\newcommand{\D}{\mathcal{D}}
\newcommand{\A}{\mathcal{A}}
\newcommand{\X}{\mathcal{X}}
\newcommand{\Z}{\mathcal{Z}}
\newcommand{\I}{\mathcal{I}}
\newcommand{\B}{\mathcal{B}}
\newcommand{\Cost}{\operatorname{Cost}}
\newcommand{\Gain}{\operatorname{Gain}}
\newcommand{\VOI}{\operatorname{VOI}}
\newcommand{\TODO}{\textbf{[TODO]}}

\title{开放世界长程 Agent 研究基础设施：\\检索、状态/时机与历史处理的统一调研（第四轮：主问题收敛与正式结论）}
\author{Research review draft}
\date{2026-09-23}

\begin{document}
\maketitle

\begin{abstract}
本文针对开放世界、长程、工具增强、多 Agent/Swarm 的算法研究系统，系统拆解三个相互强化的核心困难：\textbf{(1) 检索困难}，尤其是跨 Git 历史、代码、实验、论文与互联网的语义检索；\textbf{(2) 时机困难}，即对当前研究状态的结构化估计以及对不同未来研究动作的价值预测；\textbf{(3) 历史处理困难}，即不断增长的 agent trajectory、tool calls、subagents、swarm、实验、成果和 Git 演化造成的 provenance/semantic/branch/policy entropy。本文将三者统一为一个“provenance-grounded, partially observed, resource-bounded adaptive control”问题，并把现有工程手段分为三种 enforcement 层：machine-enforceable code/policy、agent-executed SOP、human expert judgment。调研覆盖软件仓库挖掘与 feature location、provenance 与 scientific workflow、event sourcing、process mining/object-centric process mining、agent memory、long-horizon agents、active retrieval、rational metareasoning/value of information、predictive process monitoring、learning to defer 与 tacit expertise。核心工程结论是：不应先构造庞大的“Research OS”或固定本体，而应以 append-only event/provenance substrate 为底座，以可重建的 materialized views 作为 memory/state/index，再逐层研究 state estimation、forecasting、retrieval planning 与 human deferral。
\end{abstract}

\tableofcontents

\section{问题定义：三个困难实际上构成一个闭环}

设研究系统的不可直接观测真实状态为
\begin{equation}
 z_t\in\Z,
\end{equation}
其中可包含当前知识、候选机制、代码状态、实验结果、失败模式、未决假设、外部证据覆盖、预算与风险等。系统直接看到的是异构观测
\begin{equation}
 o_t=(\text{git},\text{files},\text{agent traces},\text{tool calls},\text{experiments},\text{papers},\text{scores},\ldots).
\end{equation}
历史为
\begin{equation}
 H_t=(o_{0:t},a_{0:t-1}).
\end{equation}
研究动作可以包含
\begin{equation}
 a_t\in\{\text{Think},\text{Experiment},\text{SearchHistory},\text{SearchCode},\text{SearchWorld},\text{Branch},\text{Merge},\text{AskHuman},\text{Stop},\ldots\}.
\end{equation}

因此三个困难可写为：
\begin{align}
\textbf{History/State:}\quad & H_t\longrightarrow b_t(z)=\Pp(z_t=z\mid H_t),\\
\textbf{Forecast/Timing:}\quad & (b_t,a)\longrightarrow \Pp(z_{t+\Delta},Y\mid b_t,a),\\
\textbf{Retrieval:}\quad & (b_t,q,B_t)\longrightarrow \text{evidence/action allocation}.
\end{align}
这与 POMDP 中 belief state、state information acquisition 与 optimal stopping 的经典结构相容\parencite{monahan1982pomdp,lauri2022pomdp}，但研究过程的状态空间、动作空间和观测结构远比标准有限 POMDP 更开放、更非平稳。

本文的统一观点是
\begin{equation}
\boxed{\text{Research process}=\text{partially observed}+\text{provenance grounded}+\text{resource bounded adaptive control}.}
\end{equation}

三个困难形成正反馈：
\begin{equation}
 H_{\rm history}\uparrow\Rightarrow
 \epsilon_{\rm state}\uparrow\Rightarrow
 \epsilon_{\rm query}\uparrow\Rightarrow
 C_{\rm retrieval}\uparrow\Rightarrow
 H_{\rm history}\uparrow.
\end{equation}
这里的 $H_{\rm history}$ 不是严格 Shannon entropy，而是后文定义的工程熵的合称。

\section{研究地图：哪些成熟领域已经覆盖了哪些子问题}

\begin{longtable}{p{0.18\textwidth}p{0.31\textwidth}p{0.42\textwidth}}
\toprule
领域 & 已有成熟问题 & 对本项目的直接意义\\
\midrule
Software repository mining / feature location & 从 source、history、issues、changesets、dynamic traces 中定位功能与演化 & 复杂 Git 仓库中的“某概念在哪里、何时改变、为什么改变”并非新问题；单一 lexical/semantic search 不足\parencite{dit2013feature,kagdi2007msr,chochlov2017historical}\\
Project memory & 将开发工件自动组织为团队项目记忆 & Hipikat 已把 issue、code、history 等视为项目群体记忆\parencite{cubranic2005hipikat}\\
Provenance & Entity--Activity--Agent、derivation、usage、attribution & 可作为最小 provenance 语义骨架；不必自造底层 ontology\parencite{w3cprov2013}\\
Scientific workflow provenance & workflow run、input/output/code、可重现打包 & Workflow Run RO-Crate 已提供跨 workflow system 的运行 provenance 表示\parencite{leo2024wfrun}\\
Event sourcing & append-only event log + replay + materialized view & “历史不可变，memory/state 是投影视图”的直接工程先例\parencite{fowler2005eventsourcing,santos2026esaa}\\
Process mining & event log $\to$ process discovery/conformance/performance & 可用于从 agent 轨迹发现实际研究过程，而非只依赖 SOP 中声明的过程\parencite{vanderAalst2016processmining,diba2021eventabstraction}\\
Object-centric process mining & 一个 event 同时关联多个对象/案例 & 非常适合 agent/\allowbreak candidate/\allowbreak hypothesis/\allowbreak experiment/\allowbreak branch 多重归属\parencite{berti2023ocpm,vanderAalst2026ocpm}\\
Predictive process monitoring & 从 running case 预测 outcome/next event/remaining time & 为“当前研究过程未来会怎样”提供直接方法论邻居\parencite{francescomarino2024ppm,verenich2019remaining}\\
Agent memory & storage/reflection/experience、state management、source-grounded evolution & 说明 trajectory preservation、compression、experience abstraction 是独立问题；纯 semantic RAG 不足\parencite{zhang2025memorysurvey,luo2026memorysurvey,chen2026mage,sun2026trajwiki}\\
Long-horizon agents & cross-context state persistence、harness、recovery & 强调长程能力是系统属性而非单纯 context window\parencite{dong2026longhorizon,wu2026reliable}\\
Active retrieval & when/what to retrieve & 已在 RAG/reasoning 中出现按需检索与 intervention policy，故“何时检索”本身不是 novelty\parencite{jiang2023flare,asai2024selfrag,guo2026realm,feng2026camdf}\\
Rational metareasoning / VOI & 对 computation/information acquisition 定价 & 为“继续内部研究还是检索/实验/问人”提供规范性框架\parencite{russell1991metareasoning,hay2012selecting,saar2009activefeature,desabbata2024rational}\\
Learning to defer & 何时把决策交给昂贵的人类专家 & 将 human-in-the-loop 从固定审查者改写为可学习的 meta-action\parencite{hemmer2023defer}\\
Tacit expertise / NDM & 经验形成 pattern/tacit knowledge 与快速判断 & 为“taste”提供比“不可形式化直觉”更可操作的解释\parencite{patterson2010tacit,klein2015intuition,polanyi1958personal}\\
\bottomrule
\end{longtable}

\section{困难一：检索不是单一 RAG，而是跨时间、结构和 provenance 的证据分配}

\subsection{大型 Git 仓库为什么难检索}

传统代码搜索常把问题近似为
\begin{equation}
 q\longrightarrow \operatorname{rank}(\text{files/symbols}).
\end{equation}
但真实维护问题往往是：
\begin{quote}
“这个不断演化的功能什么时候从 A 设计变成 B？哪些 commit、issue、实验和决策共同导致了它？当前残留约束是什么？”
\end{quote}
因此目标不是静态 source retrieval，而是
\begin{equation}
 q\longrightarrow \operatorname{rank}(\text{concepts},\text{artifacts},\text{events},\text{commits},\text{rationales},\text{experiments}).
\end{equation}

Feature location 已有数十年研究。系统综述将方法覆盖 textual/IR、static structural、dynamic 与 hybrid 等路线\parencite{dit2013feature}。历史文本也有独立价值：Chochlov 等在 8 个系统、600 个 features 上发现，changeset history 与当前源码 lexical information 的组合优于各自单独使用，但效果依赖系统和配置\parencite{chochlov2017historical}。这直接反对“一个 embedding index 解决仓库历史”的设想。

更接近我们目标的是 project memory。Hipikat 将开发过程中自然产生的工件组织成可推荐的 group memory\parencite{cubranic2005hipikat}；Codebook 则用 typed graph 统一 people、changesets、work items、files、source code 及其关系，说明“多工件图 + 路径查询”在软件工程中已有先例\parencite{begel2010codebook}。

\subsection{检索应当分层}

建议将检索动作拆为
\begin{equation}
 R(q)=\operatorname{Route}(q)\to (R_0,R_1,R_2,R_3,R_4,R_5),
\end{equation}
其中：
\begin{align*}
R_0 &: \text{exact lookup: ID/hash/path/commit/experiment};\\
R_1 &: \text{structural: dependency/provenance/parent/descendant};\\
R_2 &: \text{temporal: evolution/change point/history};\\
R_3 &: \text{lexical/semantic: BM25/embedding/code search};\\
R_4 &: \text{derived experience: similar failures/mechanisms/trajectories};\\
R_5 &: \text{external: GitHub/papers/web/packages}.
\end{align*}

检索的目标不是 recall 最大化，而是有限预算下的 evidence allocation：
\begin{equation}
 r^*=\arg\max_{r\in\mathcal R}\frac{\E[\Delta U\mid r,H_t]}{\E[\Cost(r)\mid H_t]}.
\end{equation}
这使 repository retrieval 与 open-world retrieval 统一进入 metareasoning，而不是两个独立工具。

\subsection{历史搜索的关键缺口：syntactic evolution 与 semantic evolution}

Git 原生记录的是 syntactic evolution：blob、tree、commit、diff。研究者真正需要的是 semantic evolution：
\begin{equation}
 \text{Concept}\leftrightarrow\text{Decision}\leftrightarrow\text{Artifact}\leftrightarrow\text{Commit}\leftrightarrow\text{Experiment}.
\end{equation}

因此最值得构建的不是“更聪明的 git log”，而是 temporal provenance 上的 semantic index。其最低要求是：任何语义 claim 都能回到具体 commit/diff/event/evidence；任何 commit 也能向上关联其研究意图和验证结果。

\section{困难三优先：历史处理是状态估计与检索的共同底座}

\subsection{不要把 history 与 memory 混为一谈}

近期 agent memory 文献已经从简单 storage 走向 reflection 与 experience abstraction\parencite{luo2026memorysurvey}。MAGE 更明确指出，按 semantic similarity 组织 memory 会打碎 execution-state dependencies，并提出 hierarchical state tree 以保持 active path、隔离错误分支和支持 revise\parencite{chen2026mage}。TrajWiki 则使用 source-grounded evolution trajectory、immutable episodic snapshots 与 ADD/REVISE/DEPRECATE 操作保存 memory 的演化过程\parencite{sun2026trajwiki}。

这些工作支持一个保守原则：
\begin{equation}
\boxed{\text{History is evidence; memory is a derived view.}}
\end{equation}
即不以 summary 覆盖 raw history，而是
\begin{equation}
 H_{raw}\xrightarrow{f_q}V_q,
\end{equation}
允许针对不同查询重建不同 view。

\subsection{Event sourcing 是很合适的工程原型}

Event Sourcing 的核心是将状态变化保存为不可变事件，再通过 replay/projector 形成 materialized state\parencite{fowler2005eventsourcing}。2026 年 ESAA 已直接把该模式用于 LLM software agents：agent 产生结构化 intention，deterministic orchestrator 验证并写入 append-only log，再投影为可验证 view\parencite{santos2026esaa}。

对研究基础设施，建议最小事件结构为
\begin{equation}
 e_i=(id,time,actor,type,inputs,outputs,parents,refs,cost,status,env).
\end{equation}
actor 可以是 human/agent/subagent/tool/evaluator；outputs/refs 可指向 commit、file、dataset、candidate、experiment、paper、search result、claim。

\subsection{三层历史表示}

\paragraph{$L_0$: immutable evidence.}
完整保存 tool calls、prompts/outputs（受隐私与成本策略约束）、Git objects、实验 artifact、evaluator outcome、外部来源快照/标识。

\paragraph{$L_1$: structured events.}
将异构日志归一为机器可查询事件，例如 CandidateEvaluated、ArtifactRetrieved、DecisionMade、AgentSpawned、BranchRejected。

\paragraph{$L_2$: semantic research views.}
由 $L_0/L_1$ 派生 Hypothesis、Mechanism、FailureMode、Decision、OpenQuestion、Evidence 等研究语义对象。关键约束为
\begin{equation}
 L_2\to L_1\to L_0
\end{equation}
始终可追溯。

W3C PROV 的 Entity--Activity--Agent 及 used/wasGeneratedBy/wasDerivedFrom/wasAssociatedWith 等关系足以作为 $L_1$ provenance 的最小标准骨架\parencite{w3cprov2013}。OpenLineage 的 Run/Job/Dataset + extensible facets 展示了工业级 event schema 如何保持小核心与可扩展元数据\parencite{openlineage}.Workflow Run RO-Crate 则展示了科研计算中将 inputs、outputs、code 和 run provenance 一起打包的成熟做法\parencite{leo2024wfrun}。

\subsection{为什么 object-centric process mining 很匹配 swarm/agent research}

传统 process mining 假定一个 case ID，但研究事件往往同时属于：
\begin{equation}
 \text{AgentRun},\text{Candidate},\text{Hypothesis},\text{Experiment},\text{GitBranch},\text{Task}.
\end{equation}
Object-centric process mining 允许事件同时关联多个 co-evolving objects，正是为突破单 case 假设而发展\parencite{berti2023ocpm,vanderAalst2026ocpm}。因此 swarm/subagent 不必另造历史系统；只需把 agent run 作为一种 object，并用 spawnedBy/used/generated/informed 等边描述关系。

\subsection{Event abstraction 是承重难题，不是格式问题}

底层可能记录数百万低层事件，但研究状态需要高层事件：
\begin{equation}
 \{\text{read},\text{edit},\text{test},\text{commit}\}\Rightarrow\text{ImplementMechanism}.
\end{equation}
Process mining 已把 event abstraction 视为独立研究问题，因为不同系统产生的事件粒度天然不同\parencite{diba2021eventabstraction}。因此应避免把 LLM 生成的高层摘要直接当事实；更稳妥的是把它视为可版本化、可失效、带 provenance 的 derived interpretation。

\section{困难二：时机 = 状态估计 + 未来预测 + 元决策}

\subsection{状态估计先于 policy learning}

若历史表示为 $H_t$，真正供 policy 使用的应是 belief/state representation
\begin{equation}
 b_t=\phi(H_t).
\end{equation}
若 $\phi$ 丢失关键因果结构，则即使 policy optimizer 很强也无济于事。因此建议研究顺序为
\begin{equation}
\boxed{Observability\to Representation\to Prediction\to Control.}
\end{equation}

\subsection{Predictive process monitoring 提供了直接邻居}

Predictive Process Monitoring (PPM) 从 historical event logs 学习对 running cases 的预测，包括 outcome、next activity、remaining time 与 cost 等\parencite{francescomarino2024ppm,verenich2019remaining}。对研究 agent，可以类比定义：
\begin{align}
\hat p_{\rm success} &= \Pp(\text{verified improvement within }B\mid H_t,a),\\
\hat T_{\rm gain} &= \E[\text{cost to next meaningful improvement}\mid H_t,a],\\
\hat \Delta J &= \E[J_{t+\Delta}-J_t\mid H_t,a],\\
\hat I &= \E[\text{information gain}\mid H_t,a].
\end{align}
这比只做 plateau detection 更一般。

\subsection{内部研究生产率的最小状态}

若 evaluator 产生 best-so-far $J_t^*$，可将有效 improvement 视为 marked point process。设事件到达率为 $\lambda_t$，条件改进幅度为 $\mu_t$，则
\begin{equation}
 r_I(t)\approx\lambda_t\mu_t
\end{equation}
是 empirical endogenous research productivity。其优点是区分“越来越难找到改进”（$\lambda\downarrow$）和“仍频繁改进但幅度变小”（$\mu\downarrow$）。第一代 timing policy 可以只用
\begin{equation}
 s_t=(\hat\lambda_t,\hat\mu_t,b_t),
\end{equation}
全部由 harness-authoritative evaluator trace 计算，避免 agent self-report、embedding 或 LLM judge 进入最初因果链。

\subsection{Rational metareasoning：研究动作本身也需要定价}

经典 metareasoning 将 computation 视为有成本的 meta-action，选择能提高最终决策效用的 computation\parencite{russell1991metareasoning,hay2012selecting}。Active feature-value acquisition 同样根据信息价值决定是否支付成本获取额外特征\parencite{saar2009activefeature}。这为研究 agent 给出统一形式：
\begin{equation}
 Q(a\mid b_t)=\E[U_{future}\mid b_t,a]-\Cost(a),
\end{equation}
\begin{equation}
 a_t^*=\arg\max_{a\in\A_t}Q(a\mid b_t).
\end{equation}
其中 $a$ 可以是继续内部 evolution、运行实验、检索本地历史、检索外部世界、问人或停止。

\subsection{Active retrieval 说明“何时检索”已有局部解，但研究时域不同}

FLARE 通过未来句子的低置信 token 决定 when/what to retrieve\parencite{jiang2023flare}；Self-RAG 学习按需检索和 reflection\parencite{asai2024selfrag}；ReaLM-Retrieve 把 reasoning step 上的 retrieval intervention 写成 learned sequential policy\parencite{guo2026realm}；CAM-DF 则研究异质工具成本下何时停止继续获取工具\parencite{feng2026camdf}。因此本项目的 novelty 不应表述为“首次研究检索时机”，而应收窄为：
\begin{equation}
\boxed{\begin{gathered}
\text{long-horizon algorithm research}+\text{verified-progress state}\\
+\text{history/provenance}+\text{costed open-world acquisition}.
\end{gathered}}
\end{equation}

DeepEvolve 已表明外部 deep research 与 algorithm evolution 可以组合，并以纯内部 evolution 在复杂任务中 plateau 为动机\parencite{liu2025deepevolve}。这进一步说明“外部知识是否有价值”不是最值得消耗研究资源的问题；更尖锐的问题是 acquisition policy。

\section{困难一与二的耦合：检索本身也是 sequential decision process}

检索并非一次函数调用，而可以表示为
\begin{equation}
 a_t^{ret}\in\{\text{search exact},\text{expand lineage},\text{search history},\text{semantic search},\text{search web},\text{read},\text{adapt},\text{stop retrieval}\}.
\end{equation}
一个好的 retrieval planner 必须同时预测：
\begin{enumerate}
\item 该来源命中相关证据的概率；
\item 证据的新颖性与可信度；
\item 阅读/理解/adaptation 成本；
\item 证据对当前候选 action ranking 的 expected value；
\item 继续检索相对立即行动的 opportunity cost。
\end{enumerate}

因此可定义
\begin{equation}
 \VOI(r\mid b_t)=\E[U^*_{after\ r}-U^*_{now}\mid b_t]-\Cost(r),
\end{equation}
并把 retrieval stopping 写为
\begin{equation}
 \max_r\VOI(r\mid b_t)\le 0\Rightarrow \text{StopRetrieval}.
\end{equation}
这与 cost-aware tool acquisition 和 rational metareasoning 的结构一致\parencite{feng2026camdf,russell1991metareasoning}。

\section{“系统熵”的可操作分解}

为了避免把“熵”只当比喻，建议至少记录四种工程不确定性：
\begin{align}
H_{prov} &: \text{provenance ambiguity：不知道结果/决策来自哪里};\\
H_{sem} &: \text{semantic aliasing：同一概念多名、多概念同名};\\
H_{branch} &: \text{branch explosion：agent/candidate/experiment 分支增长};\\
H_{policy} &: \text{policy conflict：code/config/MD/prompt 对行为规定冲突}.
\end{align}
它们不必强行合成一个 Shannon 标量。更实用的是为每一项定义可测 proxy，例如 unresolved provenance links、concept alias collision、active branch count、conflicting policy assertions。核心因果假设是
\begin{equation}
(H_{prov},H_{sem},H_{branch},H_{policy})\uparrow
\Rightarrow C_{retrieve}\uparrow,\ \epsilon_{state}\uparrow.
\end{equation}
该假设应通过真实 trajectory 数据验证，而非写成既定事实。

\section{三层治理：Code/IaC、SOP 与 Human Taste}

\subsection{Layer 1：machine-enforceable invariants}

IaC 的核心价值是用 machine-readable definition 取代手工配置\parencite{rahman2021iac}；Policy as Code 进一步把安全、资源、合规等组织/技术规则写成 machine-enforceable declarative policy\parencite{serra2026policy}. 对研究 agent，凡可机械验证的 invariant 均应尽量从自然语言提升为 code：
\begin{equation}
\boxed{\text{Mechanically verifiable}\Rightarrow\text{do not leave it only in prose}.}
\end{equation}
典型对象包括 schema、hash、budget、sandbox、permission、evaluator identity、artifact dependency、allowed transition、provenance completeness。

\subsection{Layer 2：SOP 是软 policy，不是事实}

SOP 更接近
\begin{equation}
 \pi_{SOP}(s)\to a
\end{equation}
而不是 invariant。例如“连续三次相似失败后先查历史”不能保证永远正确，也难以在所有环境强制。因此 SOP 应带 version、scope、applicability condition 与 evidence，而不只是散落 Markdown。Agent 对 SOP 的遵循率、偏离原因和 outcome 应可测。

\subsection{Layer 3：Human taste = 尚未充分显式化的 long-horizon evaluator}

关于 expert intuition，更谨慎的表述不是“人类有神秘 taste”，而是经验可形成 pattern recognition 与 tacit knowledge；Naturalistic Decision Making 文献把熟练直觉与经验形成的 pattern/tacit knowledge 联系起来\parencite{patterson2010tacit,klein2015intuition}。Polanyi 的 personal/tacit knowledge 则提醒我们，科学实践包含无法完全还原为显式规则的个人知识成分\parencite{polanyi1958personal}。

工程上可将 taste 暂时建模为
\begin{equation}
 T_H(s)\approx \E[\text{long-horizon research value}\mid s,\text{human experience}],
\end{equation}
但必须强调：这是计算解释，不是已被心理学证明的等价关系。

\subsection{Human override 应成为一等事件}

记录
\begin{equation}
 h_i=(state,agentProposal,humanDecision,reason,confidence,outcome).
\end{equation}
这使 human-in-the-loop 从不可分析的聊天干预变成可学习数据。Learning to defer 已研究如何在专家预测昂贵时学习何时交给人类\parencite{hemmer2023defer}。对研究 agent，未来可学习
\begin{equation}
 \pi_{defer}(b_t)=\Pp(\text{human intervention has positive net value}\mid b_t).
\end{equation}

\subsection{知识结晶：Taste $\leftrightarrow$ SOP $\leftrightarrow$ Code}

成熟系统不应固定三层边界，而应支持知识在 enforcement level 间迁移：
\begin{equation}
 \text{Taste}\xrightarrow{codify}\text{SOP}\xrightarrow{mechanize}\text{Invariant}.
\end{equation}
若 hard invariant 出现系统性 counterexample，也允许
\begin{equation}
 \text{Invariant}\xrightarrow{invalidate/revise}\text{SOP/Taste}.
\end{equation}
关键是每次升级/降级都有 evidence 与 provenance，而不是“文档悄悄变了”。

\section{建议的最小体系结构：Control Plane 与 Epistemic Plane 分离}

\subsection{Control Plane}
回答：什么可以发生、必须发生、不得发生。
\begin{itemize}
\item schemas/contracts;
\item sandbox/network/filesystem policy;
\item evaluator identity and budget;
\item artifact integrity;
\item executable transition guards;
\item reproducibility/freeze checks.
\end{itemize}

\subsection{Epistemic Plane}
回答：发生了什么、当前相信什么、为什么、还不知道什么、下一步去哪找。
\begin{itemize}
\item append-only research event ledger;
\item provenance graph;
\item semantic/temporal indexes;
\item state estimators;
\item predictive models;
\item retrieval planner;
\item human override/defer records;
\item derived research views.
\end{itemize}

两者通过 stable IDs 与 content hashes 连接，但不互相替代。

\section{Event-first, schema-light, views-later}

本文不建议立即设计完整 Research Ontology。理由有三：
\begin{enumerate}
\item 研究对象和 agent 行为仍在快速变化，过早本体化会冻结错误抽象；
\item object-centric process management 明确指出不同目的需要不同 process representation，不存在单一模型满足所有用途\parencite{vanderAalst2026ocpm};
\item agent memory 文献也正在从静态 record 转向 source-grounded、可演化 memory，而非证明某个固定 schema 是最终答案\parencite{sun2026trajwiki,chen2026mage}。
\end{enumerate}

因此建议只冻结极小 core：
\begin{verbatim}
Event {
  id, timestamp, actor, type,
  inputs[], outputs[], parents[], refs[],
  cost, status, environment_id, schema_version
}
\end{verbatim}
所有 Hypothesis/Mechanism/Decision/Failure 等均先作为 typed derived objects/views，而不是底层不可变 ontology。

\section{建议的研究程序：不要一次性解决三个困难}

\subsection{Phase 0：Observability audit}
目标：回答“我们到底记录了什么，哪些东西丢了？”
\begin{itemize}
\item 对 agent/tool/git/evaluator/subagent/swarm 事件做 source inventory；
\item 定义 authoritative event IDs 与 clocks；
\item 测量 provenance completeness；
\item 识别不可恢复的 silent state mutation。
\end{itemize}

\subsection{Phase 1：Append-only substrate}
先实现事件 ledger + content-addressed artifact references + Git linkage，不做复杂 memory AI。
验收：任一最终 result 可反向恢复其 candidate、experiment、agent run、tool inputs 与 commit；任一 event 可 replay/verify 到约定边界。

\subsection{Phase 2：Retrieval baselines}
建立五类 baseline：exact、graph/provenance、temporal Git/history、lexical/semantic、hybrid。用真实历史问题构造 benchmark，例如：
\begin{itemize}
\item 某机制首次出现在哪里？
\item 为什么被废弃？
\item 当前实现继承了哪些旧假设？
\item 哪些实验已经证伪相似方案？
\item 某功能跨 rename/move/refactor 后的完整 lineage 是什么？
\end{itemize}
指标除 Recall@k/MRR 外，应测 evidence completeness、time-to-answer、wrong-rationale rate、provenance precision。

\subsection{Phase 3：Event abstraction / state views}
比较 deterministic projection、rule-based abstraction、LLM abstraction、hybrid。高层 view 只做 derived object，并保留 source links。重点测：
\begin{equation}
\text{compression}\quad vs\quad\text{future-query recoverability}.
\end{equation}

\subsection{Phase 4：State estimation}
先不用 RL。用历史 campaign 学习/校准：
\begin{equation}
 \hat\lambda_t,\hat\mu_t,\hat p_{success},\hat T_{gain},\hat p_{failure}.
\end{equation}
评价 calibration、Brier score、time-dependent discrimination 与 distribution shift。

\subsection{Phase 5：Action forecasting}
针对固定 action set 比较
\begin{equation}
 \hat Q(a\mid b_t)
\end{equation}
的预测质量。先离线 counterfactual/paired experiments，再考虑 online policy learning。

\subsection{Phase 6：Retrieval timing}
使用受控 exogenous packet 做因果 timing assay，再进入 live open-world search。这样分离 information timing 与 query formulation/source quality。

\subsection{Phase 7：Human deferral}
把 human intervention 变成显式动作，比较 always-human、never-human、uncertainty threshold、learned defer policy。

\subsection{Phase 8：Knowledge crystallization}
从高频 human override / SOP violation / repeated failure 中挖掘候选 rule，经过 replay 和 counterexample testing 后再提升为 executable policy；禁止 LLM 自动把“常见经验”直接升级成 hard invariant。

\section{实验设计中的关键反例与风险}

\subsection{Summary trap}
摘要在当前 query 上有效，不代表对未来 query 足够。故 raw evidence 不应被 summary 替代。

\subsection{Graph trap}
有 provenance graph 不等于能回答语义问题；图的 edge typing 和 query planner 可能成为新瓶颈。

\subsection{Ontology trap}
过早固定 Hypothesis/Mechanism/Decision ontology 会迫使新型研究行为映射到旧类型，制造 silent semantic loss。

\subsection{Embedding trap}
相似不等于因果/适用。Agent memory 实证研究显示 retrieved experience 会强烈诱导 experience-following behavior，因此错误或过时 memory 可能系统性污染后续动作\parencite{xiong2026experience}。

\subsection{Prediction trap}
若训练数据来自旧 policy，则 $P(Y\mid s,a)$ 存在 policy-induced selection bias；需要 exploration、off-policy correction 或随机化 intervention 才能逐渐获得 action-value 识别。

\subsection{Human oracle trap}
Human taste 不是 ground truth。应记录 human confidence、分歧、结果和反事实；人类 override 也要被 calibration。

\subsection{Automation ratchet trap}
Taste$\to$SOP$\to$hard rule 是危险的单向棘轮。必须允许 hard rule 被 counterexample 降级或失效。

\section{理论上最值得研究的开放问题}

\subsection{Query-conditioned sufficient history}
给定未来 query $q$，寻找最小历史投影
\begin{equation}
 V_q=f_q(H)
\end{equation}
使得任务损失增加不超过 $\epsilon$：
\begin{equation}
 L(q,V_q)-L(q,H)\le\epsilon,
\end{equation}
同时最小化存储/读取成本。这是 memory compression、retrieval 与 provenance 的统一目标。

\subsection{Research-state representation learning}
寻找表示 $z_t=\phi(H_t)$，使其同时对未来 outcome/action value 有预测力，并保持 provenance 可解释性。纯 embedding 与纯手工 schema 都可能不足。

\subsection{Action-value under open-world expansion}
动作集本身会随检索发现新工具、算法、数据而扩张：
\begin{equation}
 \A_{t+1}=\A_t\cup\Delta\A_t.
\end{equation}
这使标准固定动作空间的 POMDP/RL 假设失效。需要把 SearchWorld 同时视为信息获取和 action-space expansion。

\subsection{Attribution across long trajectories}
最终 improvement 往往由很早的 retrieval、某 subagent 的失败、后续 adaptation 共同产生。近期 trajectory attribution 工作开始形式化 long-range attribution，但该问题仍很开放\parencite{chen2026attribution}。

\subsection{When to reuse vs when to re-derive}
长 trajectory 检索后如何适配到当前 state 是独立瓶颈；近期工作已显示“找到相似 trajectory”与“安全有效复用”不是同一个问题\parencite{li2026qcr}。因此 retrieval quality 与 reuse quality 必须分开评测。

\subsection{Human taste distillation without Goodharting}
若把 human override 学成规则，系统会逐渐优化人类历史判断的 proxy，而非真实研究价值。需要保留 outcome-based validation 与 disagreement sampling，避免把 taste 静态化。

\section{工程建议：现在应该复用什么，不该自造什么}

\subsection{优先复用}
\begin{itemize}
\item provenance vocabulary：W3C PROV；
\item research packaging：RO-Crate / Workflow Run RO-Crate；
\item run lineage/event API 思想：OpenLineage；
\item event-sourcing pattern：append-only log + projector + replay；
\item object-centric event/process ideas：OCEL/OCPM；
\item repository mining：Git native history + symbol/AST/dependency + lexical/IR hybrid；
\item policy enforcement：IaC/Policy-as-Code 风格的 machine-enforceable contract。
\end{itemize}

\subsection{暂不自造}
\begin{itemize}
\item 完整 Research Ontology；
\item 单一“万能 memory”；
\item 单一 embedding/vector DB 作为 source of truth；
\item 未有可识别 state 前的复杂 RL timing policy；
\item 把 LLM summary 当 authoritative history；
\item 自动将 human/SOP 经验提升为 hard rule。
\end{itemize}

\section{一个可落地的最小数据模型}

\begin{verbatim}
Event
  id: content/random stable id
  ts: logical + wall clock
  actor_id: Human | AgentRun | Tool | Evaluator
  type: extensible string
  input_refs[]
  output_refs[]
  parent_event_ids[]
  object_refs[]      # task, candidate, hypothesis, branch, experiment...
  policy_refs[]      # code/SOP/prompt versions that governed action
  env_id
  cost
  status
  schema_version

Artifact
  id: content hash where possible
  media_type
  locator
  git_commit?
  external_uri?
  license?
  created_by_event

Interpretation (derived, versioned)
  id
  type: Claim | Hypothesis | Mechanism | FailureMode | Decision | ...
  statement
  source_event_ids[]
  confidence
  valid_from / deprecated_by
  interpreter_version
\end{verbatim}

这个 schema 的关键不是字段本身，而是：Event/Artifact 尽量事实化；Interpretation 明确是 derived、可修订、可废弃。

\section{总判断}

这次调研后的核心判断有五点。

第一，三个困难确实都有庞大的既有研究，但它们分散在 software evolution、process mining、provenance、agent memory、RAG/metareasoning 与 human-AI decision making 中；尚没有一个成熟框架直接覆盖“开放世界长程算法研究 agent”的完整闭环。

第二，\textbf{历史处理是更底层的承重项}。没有 provenance-grounded history，就无法可靠构造 state；没有 state，就很难研究 timing；没有语义/时间结构，就会不断把历史检索退化成昂贵的重新阅读。

第三，最稳健的底座不是“Research OS”，而是
\begin{equation}
\boxed{\text{append-only event/provenance substrate}+\text{rebuildable views}.}
\end{equation}
这与 W3C PROV、Event Sourcing、OpenLineage、Workflow Run RO-Crate 和 OCPM 的成熟思想一致。

第四，IaC/SOP/Human Taste 不是互斥方案，而是不同 uncertainty/enforcement levels；成熟系统应允许知识在三层间以有证据的方式结晶和回退。

第五，真正的新研究问题位于这些成熟基础设施之上：
\begin{equation}
\boxed{
H_t\xrightarrow{\phi}b_t
\xrightarrow{\text{forecast}}\hat Q(a\mid b_t)
\xrightarrow{\text{meta-policy}}a_t,
}
\end{equation}
其中动作同时包括内部研究、历史/外部检索、实验、分支、停止与 defer-to-human；而每个 state/forecast/decision 又必须能够沿 provenance 回到原始证据。

换言之，最值得长期研究的统一问题不是“怎样给 agent 加 memory”，而是：
\begin{quote}
\textbf{如何从不断增长、异构、多主体、可分叉的研究历史中，形成可追溯的当前 belief state；如何预测不同研究动作的未来价值；以及如何在有限预算下决定何时检索、实验、分叉、停止或交给人类。}
\end{quote}


\part{第二轮深挖：从“日志”到“可控制的研究认知系统”}

\section{第二轮检索带来的结构性修正}

第一轮把总问题概括为 provenance-grounded、resource-bounded adaptive control。第二轮沿 software evolution、agent observability、object-centric process mining、predictive/prescriptive monitoring、metareasoning、agent memory evolution 与 expert judgment 分别下钻后，得到一个更严格的分层：
\begin{equation}
\boxed{
\text{Evidence Plane}
\to \text{State Plane}
\to \text{Forecast Plane}
\to \text{Control Plane}
\to \text{Action}
}
\end{equation}
并存在贯穿全链的 \emph{Judgment Plane}。这里 Evidence Plane 不是普通日志；它必须同时支持 execution telemetry、artifact lineage、semantic history 与证据追溯。OpenTelemetry 的 GenAI semantic conventions 已经开始标准化 agent/model/tool 的 trace、metrics 与 events；2026 年的 agent provenance 综述也把 evidence tracing、tool provenance、memory lineage、observability、debugging 与 recovery 统一到了 execution provenance 视角\cite{otel2026genai,wang2026provenance}。这说明“把 agent 行为完整记下来”正在从框架私有功能变成独立基础设施问题。

但是 telemetry 仍然只回答“发生了什么”。研究系统还必须回答：某次调用改变了哪个假设？哪个 commit 实现了哪个机制？哪个负结果使哪个方向被废弃？因此我们需要把传统 span tree 扩展为 research lineage graph：
\begin{equation}
\boxed{
\text{Trace} \subset \text{Provenance} \subset \text{Research History}.
}
\end{equation}

\section{困难一再分解：检索不是一个 IR 问题，而是五种不同的定位问题}

\subsection{从 feature location 到 repository navigation}

大型仓库中的“找代码”已有长期研究传统。传统 feature/concept location 组合文本、静态结构、动态行为与历史；semantic history slicing 更进一步，把 commit sequence 从纯时间序列重构为满足某一语义判据的历史切片\cite{li2018semantic}. 这与我们问“某功能为何变成今天这样”高度一致：目标不是返回最相似的 diff，而是恢复对当前语义状态有因果贡献的 change set。

最新 repository-agent 工作又给出三个互补方向。AIRCoder 表明 repository retrieval 不能只用单一相似度，而应按 query 自适应融合 textual similarity、dependency 与 structural hierarchy\cite{shi2026aircoder}；RepoLens 把细粒度代码重组为 repository-wide concerns，以缓解 concern scattering/mixing\cite{wang2025repolens}；RepoNav 则显示即使已经检索到正确 snippets，平铺结果仍不足以支持 agent 定位，file-centered navigation scaffold 本身会改善 localization\cite{chai2026reponav}。因此：
\begin{equation}
\boxed{\begin{gathered}
\text{Retrieval Quality}\neq\text{Embedding Similarity};\\
\text{Retrieval}=\text{Candidate Generation}+\text{Structural Navigation}\\
+\text{Evidence Organization}.
\end{gathered}}
\end{equation}

\subsection{时间维：当前代码检索与历史检索应分开}

对于不断调整的功能，至少存在五种 query：
\begin{enumerate}
  \item \textbf{where-now}: 当前实现在哪里；
  \item \textbf{why-now}: 为什么形成当前实现；
  \item \textbf{when-changed}: 某语义何时发生改变；
  \item \textbf{what-caused}: 哪些历史变化对当前行为有语义贡献；
  \item \textbf{what-failed}: 哪些已尝试方案后来被回退/否决，以及为什么。
\end{enumerate}
前两类可由 code graph + semantic retrieval 部分解决；后三类要求 temporal-semantic index。Semantic history slicing 已证明“以测试/行为作为 slicing criterion”可以把功能相关 commits 从完整历史中剥离出来\cite{li2018semantic}。工程上也已有 Semcode 这类工具把 symbol/call graph、普通 commit search 与 semantic commit search 放在同一索引中\cite{semcode2026}。因此我们不应再把 Git 只作为 artifact store，而应将其作为 research event graph 的一个时间化子图。

\subsection{检索本身应成为可计量阶段}

2026 年 issue-localization 研究已经明确把 repository exploration 当成独立、昂贵、可预算的 agent stage，并测量早期 candidate discovery、Hit@3、F1、agent time 与 token usage\cite{awad2026exploration}。这给我们的直接启发是：
\begin{equation}
\boxed{
C_{retrieve}=C_{locate}+C_{inspect}+C_{verify}+C_{integrate}
}
\end{equation}
而不是只统计 search API 次数。一个 retrieval action 的输出也不应是文本，而应是带 provenance 的 evidence bundle：
\begin{equation}
R=(q,\mathcal C,\mathcal E,\mathcal P,\text{cost},\text{coverage},\text{uncertainty}).
\end{equation}

\section{困难三再分解：历史处理的真正对象是多对象共同演化}

\subsection{为什么线性 trajectory 天生不适合 swarm research}

传统 case-centric event log 假定一个 event 属于一个 case；但 research event 常同时关联：
\[
\text{AgentRun},\ \text{Hypothesis},\ \text{Candidate},\ \text{Experiment},\ \text{Artifact},\ \text{GitBranch},\ \text{Task}.
\]
OCEL 2.0 已经提供 event--object、object--object relation、time-varying object attributes 与 qualified relations\cite{ocel2023}; 2026 object-centric process management manifesto 更明确把“co-evolving data objects”视为传统 case silo 的根本替代\cite{seidel2026ocpm}。这几乎直接给出了我们的底层事件模型：
\begin{equation}
\boxed{
Event\xleftrightarrow{qualified}Objects,
\qquad Objects\xleftrightarrow{}Objects.
}
\end{equation}
因此 swarm/subagent 不应产生另一套 memory schema；它们只是 AgentRun objects 及其 spawn/use/generate relations。

\subsection{低层 trace 到研究事件：event abstraction 是承重问题}

原始 telemetry 的粒度通常是 token/model call/tool call/file write，而研究 reasoning 的粒度是“验证假设”“发现机制冲突”“回退路线”“确认 prior art”。Process-mining 的 event-abstraction 文献早已指出 mixed-granularity logs 会导致模型复杂、隐藏真正过程逻辑，并系统研究了从低层 event 到高层 activity 的抽象\cite{vanzelst2021abstraction}。因此 Research Event Abstraction 不应只是 LLM 总结，而应成为可单独评估的函数：
\begin{equation}
\phi:\mathcal E_{low}^{*}\to \mathcal E_{research}.
\end{equation}
至少要评估：boundary detection、type accuracy、argument/entity preservation、provenance completeness 与 downstream retrieval utility。

\subsection{不可逆 summary 应被禁止成为唯一记忆}

近期 agent memory 工作越来越支持“保留演化历史而非覆盖”。Cognitive Scaffold 把 fluid working context 与 persistent structured memory 分离\cite{cognitive2026scaffold}；TrajWiki 使用 immutable episodic snapshots 与 ADD/REVISE/DEPRECATE 维护 claim evolution\cite{sun2026trajwiki}；EvoMem 直接把动态环境变化编码为 patch-based memory history\cite{xu2026evoarena}。2026 agent-memory systems characterization 还显示 memory 的 construction、retrieval、generation 三阶段存在明显的成本搬移与 freshness/latency trade-off\cite{omri2026memorysystems}。因此我们的原则应升级为：
\begin{equation}
\boxed{
\text{Lossless Evidence} \to \text{Versioned Derived State} \to \text{Query-conditioned Views},
}
\end{equation}
而非“定期把历史压成一个更短 MD”。

\subsection{从 experience 到 procedural knowledge：历史不是只为了回忆}

ACL 2026 的综述把 agent memory 演化明确划为 Storage $\to$ Reflection $\to$ Experience\cite{luo2026storage}; procedural-memory 工作则尝试把 trajectory 转化为可复用 how-to knowledge\cite{cao2026procedural}。Memento-Skills 更工程化地把 skill 作为可 retrieve/execute/evaluate/repair/rewrite 的一等能力对象\cite{zhou2026memento}。这提示历史系统至少要支持两条不同的派生路径：
\begin{align}
H &\to \text{episodic evidence},\\
H &\to \text{procedural experience}.
\end{align}
前者回答“发生过什么”，后者回答“下次应该怎样做”。二者不能混为一个 embedding store，因为 procedural abstraction 的错误会反过来系统性污染未来 trajectory。

\section{困难二再分解：时机问题应借鉴 predictive/prescriptive process monitoring}

\subsection{State estimation 不是自然语言总结}

若研究过程被表示为 object-centric event history $H_t$，当前状态不是某段摘要，而是 belief：
\begin{equation}
 b_t(z)=P(z_t=z\mid H_t,\mathcal V_t),
\end{equation}
其中 $\mathcal V_t$ 是当前可见的 derived views。Predictive Process Monitoring (PPM) 正是从正在运行的 event prefix 预测 next activity、remaining time、outcome 等。2026 的 controlled benchmark 甚至发现，在标准 PPM 任务上，传统 sequence models 对 next-activity prediction 仍可优于更昂贵的 LLM，而 tabular foundation models 在时间任务上具有竞争力\cite{fertig2026ppm}。这对我们是重要的反 hype 证据：
\begin{equation}
\boxed{
\text{State/forecast layer should be model-class agnostic.}
}
\end{equation}
不要预设“必须用 LLM world model”。

\subsection{从预测到干预：Prescriptive Process Monitoring 是 timing 的直接邻居}

Prescriptive Process Monitoring 的问题不是“会不会失败”，而是“运行中何时触发何种 intervention 才改善 outcome”。系统综述把它分解为 objective、metrics、intervention types、modeling technique、data input 与 intervention policy，并指出 causality 与 second-order effects 是主要开放问题\cite{kubrak2022prescriptive}。Action-oriented process mining 则更明确地试图打通 diagnostics/prediction 与 action\cite{gronau2022action}. 这与我们的 timing 问题高度同构：
\begin{equation}
\text{event prefix}\to \text{risk/value forecast}\to \text{intervention trigger}.
\end{equation}
区别只在于 intervention 从业务动作换成：
\[
RetrieveExternal,\ RetrieveHistory,\ Experiment,\ Branch,\ Stop,\ AskHuman.
\]
因此 timing research 应主动吸收 PPM 的 evaluation methodology，而不是只与 RAG retrieval timing 比较。

\subsection{为什么简单 plateau trigger 可能系统性错误}

经典 costly-information optimal stopping 已经证明：在部分可观测 Markov process 中，“购买信息”的最优区域不一定是凸的，value of information 随知识状态可能非单调\cite{monahan1980}. 这直接否定一个过强直觉：
\[
\text{progress 越慢}\Rightarrow\text{越应该检索}.
\]
真实策略可能是：早期不检索；中期检索；某些后期状态反而不检索，因为剩余预算不足以吸收信息。故 B0b 的 $\hat\lambda\hat\mu$ plateau policy 应被视为 baseline，而不是理论目标。

\subsection{Rational metareasoning 提供统一控制语言}

Russell--Wefald 的核心是把 computation 本身视为有成本动作，并按其改善外部决策质量的 expected value 来选择 computation\cite{russell1991meta}。因此 research action 可以统一为 metalevel actions：
\begin{equation}
 a_t\in\{Think,SearchHistory,SearchWorld,Experiment,Branch,AskHuman,Stop\}.
\end{equation}
定义
\begin{equation}
 VOC(a\mid b_t)=\mathbb E[U\mid b_t,a]-\mathbb E[U\mid b_t]-C(a),
\end{equation}
则理想 policy 是：
\begin{equation}
\pi^*(b_t)=\arg\max_a VOC(a\mid b_t).
\end{equation}
但由于真实 $VOC$ 不可直接计算，我们的研究核心就变成学习/校准 $\widehat{VOC}$，而不是直接训练一个 opaque policy。

\section{检索时机：从 active RAG 到 long-horizon research control}

Active Retrieval 已经明确证明“是否检索”是独立问题。UAR 把多种 retrieval criteria 统一为低开销分类任务\cite{cheng2024uar}；ReaLM-Retrieve 在 reasoning-step 粒度估计 uncertainty，以决定 reasoning 过程中何时检索\cite{guo2026realm}. DeepEvolve 又已经把 deep research 与 algorithm evolution 接入同一反馈循环\cite{liu2025deepevolve}。因此我们不能声称 novelty 是“让算法 agent 检索外部信息”或“动态决定是否检索”。

更窄而仍未被这些工作解决的问题是：
\begin{equation}
\boxed{\begin{gathered}
\text{在数十分钟至数日、包含代码/实验/git/swarm 的研究 trajectory 中，}\\
\text{如何从 provenance-grounded state 估计不同 acquisition action 的长期价值？}
\end{gathered}}
\end{equation}
这里的 acquisition 不是一次 RAG call，而可能触发一串昂贵后果：定位、阅读、冻结 artifact、适配、验证、重新实验。因此 action cost 本身具有 duration 与 delayed effect。

\section{Human Taste：从“最后兜底”升级为可研究的 judgment channel}

\subsection{Taste 不应被粗暴等同为不可形式化直觉}

Naturalistic Decision Making (NDM) 将专家直觉解释为长期经验形成的 pattern 与 tacit knowledge，并强调复杂、动态、不确定环境中的 recognition 与 sensemaking\cite{klein2015intuition,patterson2010tacit}。统计学关于 expert judgment 的讨论也强调 judgment 贯穿问题提出、设计、分析与解释，而且关键决策应留下书面记录以支持透明性\cite{brownstein2019judgment}。这支持我们的工程解释：Human Taste 不是 oracle，而是一个具有 domain coverage、calibration 与 failure modes 的预测器。

\subsection{Human intervention 应变成 supervised meta-data}

每次人工介入至少记录：
\begin{equation}
 h_i=(b_i,a_i^{agent},a_i^{human},r_i,c_i,y_i,\gamma_i),
\end{equation}
其中 $r_i$ 是理由、$c_i$ 是介入成本、$y_i$ 是后验 outcome、$\gamma_i$ 是 confidence。这样可研究三个问题：
\begin{enumerate}
 \item 哪些 state 上 human override 有正的 treatment effect？
 \item 哪些 override rationale 可结晶为 SOP？
 \item 哪些 SOP 经重复验证后可提升为 hard invariant？
\end{enumerate}
Learning-to-defer 已把“模型还是专家”作为可学习分配问题，并已有公开 benchmark framework\cite{alves2025defer}。我们的不同点是 defer target 不是单次 label，而是 research-process control decision。

\subsection{知识结晶应是可逆、有证据的状态机}

因此 IaC/SOP/Taste 三层可改写为：
\begin{equation}
\text{Judgment}\rightleftarrows\text{Procedure}\rightleftarrows\text{Invariant}.
\end{equation}
升级条件不是“人类写进文档”，而是累积 evidence：
\begin{equation}
Promote(k)\iff
N_k\ge n_{min}\land
\widehat{Benefit}_k>\delta\land
\widehat{FailureRisk}_k<\epsilon.
\end{equation}
降级则由 counterexample、environment shift 或 conformance failure 触发。这样“taste $\to$ SOP $\to$ code”从口号变成可审计的 knowledge lifecycle。

\section{一个更具体的最小架构：双账本、四平面、多个视图}

第二轮调研后，我不再建议只有一个 Event Log，而建议两个逻辑账本（物理上可同库）：
\begin{description}
 \item[Execution Ledger] 记录模型调用、tool call、agent spawn、文件/网络/评估、token/cost/latency 等可观测执行事实；尽量兼容 OpenTelemetry GenAI conventions\cite{otel2026genai}。
 \item[Research Ledger] 记录 hypothesis、decision、experiment、artifact、claim、rejection、human override 及其 provenance relation；其关系模型可借 W3C PROV 与 OCEL，而 workflow-run 打包可借 Workflow Run RO-Crate\cite{w3cprov2013,ocel2023,soiland2024wfrun}。
\end{description}
二者通过稳定 IDs 连接：
\begin{equation}
ResearchActivityID\leftrightarrow TraceID\leftrightarrow ArtifactHash\leftrightarrow CommitSHA.
\end{equation}
上层四平面：
\begin{enumerate}
 \item Evidence Plane：immutable/raw evidence；
 \item State Plane：versioned derived beliefs/views；
 \item Forecast Plane：action-conditional predictions；
 \item Control Plane：hard constraints + meta-policy。
\end{enumerate}
Judgment Plane 横跨四层，但所有介入必须落账。

\section{为什么“Research OS”仍然不是正确抽象}

2026 scientific-workflow terminology review 指出，数百个 workflow systems 已存在，但由于研究过程与基础设施差异巨大，仍未出现 turnkey universal solution\cite{workflowterms2026}；最新 AI-driven scientific workflow systems review 也把 open problems 指向 dynamic workflow representation、state-aware recovery、decision provenance 与 reproducible adaptive execution\cite{alnasir2026workflows}。这反而支持我们的 shelf-first 判断：不应再造一个全包平台，而应把缺口缩到可组合的中间层。

值得注意的是，OpenAI4S 已经采用 persistent runtime + research sessions + append-only Action Ledger + versioned artifacts + checkpoints 的路线\cite{openai4s2026}；另一个 2026 agentic scientific workflow 架构则采用“LLM semantic intent + deterministic workflow generation + Markdown Skills”的三层分解，把 nondeterminism 限制在语义解释层\cite{balis2026workflow}. 它们分别验证了我们两条原则的工程合理性：append-only provenance 与 soft knowledge/hard execution 分离。

\section{新的研究问题树}

据此，总问题可拆成六个相互连接、但可独立实验的 research programs：

\subsection*{RP1: Temporal-Semantic Retrieval}
输入 query 与 repository/research history，输出最小充分 evidence subgraph。核心指标：Recall@budget、causal/semantic history slice quality、time-to-evidence、downstream decision gain。

\subsection*{RP2: Research Event Abstraction}
学习/设计 $\phi:\mathcal E_{low}^{*}\to\mathcal E_{research}$。核心指标：event boundary、type、entity preservation、provenance completeness、reconstructability。

\subsection*{RP3: Belief-State Estimation}
从 object-centric event history 构造 $b_t$。核心不是摘要可读性，而是 sufficiency：在固定预算下，$b_t$ 是否保留了预测/控制所需信息？可用 predictive sufficiency 与 retrieval regret 测试。

\subsection*{RP4: Research Foresight}
学习 $P(\Delta U,\Delta I,C,Failure\mid b_t,a)$。先比较简单统计模型、survival/hazard、sequence model、tabular FM、LLM world model，再决定复杂度；PPM 的最新结果说明不能默认 LLM 胜出\cite{fertig2026ppm}。

\subsection*{RP5: Metareasoning / Timing}
在相同 state/forecast 下比较 fixed、threshold、myopic VOC、lookahead、learned policy。必须显式包含 Stop 与 AskHuman，并评估 anytime utility、cost、regret 与 calibration。

\subsection*{RP6: Judgment Crystallization}
从 human overrides 与 agent trajectories 中发现可复用 procedure，再验证是否值得升级为 SOP/invariant。核心指标：override value、rule precision/coverage、distribution shift robustness、promotion/revocation latency。

\section{一个关键新概念：研究基础设施的目标不是降低信息熵，而是降低决策相关熵}

“系统熵”若只定义为文件多、日志多，会诱导过度压缩。真正相关的是 action uncertainty。设完整历史为 $H_t$，未来动作价值为 $Q(a)$，则我们真正希望 state representation $S_t=\phi(H_t)$ 尽量保持：
\begin{equation}
I(S_t;Q(A))\approx I(H_t;Q(A)),
\end{equation}
同时压缩成本：
\begin{equation}
|S_t|\ll |H_t|.
\end{equation}
因此更准确的目标不是最小化 $H(H_t)$，而是最小化\textbf{decision-relevant information loss}：
\begin{equation}
\mathcal L_{DR}=I(H_t;Q(A))-I(S_t;Q(A)).
\end{equation}
这给 event abstraction、memory compression、semantic indexing 一个共同的评价标准：一个 view 是否“好”，最终应由它是否保持未来研究决策所需的信息来判断，而不是由 summary 是否漂亮来判断。

\section{工程优先级：现在最值得做什么}

结合成熟度与研究价值，建议顺序改为：
\begin{enumerate}
 \item \textbf{P0 Telemetry/ID unification}：先让 agent/model/tool/evaluator/git/artifact 有稳定关联 ID；优先复用 OpenTelemetry，不自造 trace stack。
 \item \textbf{P1 Object-centric append-only ledger}：最小 schema，兼容 OCEL/PROV 思路；不先做巨大 ontology。
 \item \textbf{P2 Temporal-semantic repository index}：把 symbol/dependency/current semantic search 与 commit/history search 联合；这是最直接缓解当前痛点的 shelf-first 项。
 \item \textbf{P3 Derived research views}：Hypothesis/Experiment/Failure/Decision/Artifact views，全部可回溯。
 \item \textbf{P4 Human override ledger}：从现在开始积累 taste 数据，成本极低、未来价值高。
 \item \textbf{P5 Offline process mining}：先做 discovery/conformance/event-abstraction 分析，不立即闭环控制。
 \item \textbf{P6 Forecast benchmark}：从历史 campaign 构造 next-action/outcome/cost prediction dataset，比较简单模型到 LLM。
 \item \textbf{P7 Closed-loop timing}：只有 state/forecast 可靠后再学习 retrieval/experiment/stop/defer policy。
\end{enumerate}

\section{第二轮结论}

第二轮调研把总问题进一步从“agent memory + retrieval + timing”收敛成：
\begin{equation}
\boxed{
\textbf{Temporal-semantic research process intelligence}
}
\end{equation}
其最小闭环为：
\begin{equation}
\boxed{
\underbrace{Event/Provenance}_{\text{事实}}
\to
\underbrace{Abstraction/Belief}_{\text{现状}}
\to
\underbrace{Forecast}_{\text{未来}}
\to
\underbrace{Metareasoning}_{\text{时机}}
\to
\underbrace{Action}_{\text{研究行为}}
\to Event.
}
\end{equation}
检索困难是“从历史/世界中定位最小充分 evidence”；历史困难是“如何让多对象、多主体、多分支演化保持可追溯且可重新投影”；时机困难是“如何从 belief state 预测 action-conditioned future value 并做 resource-bounded control”。Human Taste 则不是第四个孤立模块，而是当前最强但未充分结构化的 judgment/forecast channel。

这也给出了一个非常明确的研究边界：底层 telemetry、provenance、event format、workflow packaging、基础 code search 应尽量复用成熟标准与工具；真正值得我们投入原创研究的地方，是 \emph{research-specific event abstraction、temporal-semantic retrieval、belief-state sufficiency、action-value forecasting 与 judgment crystallization}。这几项恰好对应当前文献之间尚未被打通的接口。



\part{第三轮深挖：从问题地图进入可证伪研究程序}

\section{为什么广度与深度必须交替，而不是二选一}

继续检索后，一个方法论结论变得更清楚：本问题不是单一成熟学科内的优化问题，而是多个成熟学科接口处的\textbf{开放发现问题}。若一直横向扩展，会产生文献枚举而缺少可检验结构；若过早纵向钻入一个分支，又会把局部术语误当成总问题。文献检索本身已有两个可借鉴传统。其一，systematic review 中 database search 与 backward/forward snowballing 是互补的搜索机制\parencite{wohlin2014snowballing,wohlin2015experiences}；其二，Literature-Based Discovery 区分 open discovery 与 closed discovery：前者从一个问题寻找未知连接，后者在已知两端之间寻找机制链\parencite{henry2017lbd,gopalakrishnan2019lbd}。

据此，本文建议把研究推进本身写成一个 meta-policy：
\begin{equation}
 m_k\in\{\text{Breadth},\text{Depth},\text{Synthesis},\text{Experiment}\}.
\end{equation}
Breadth 的目标是降低“未知邻接领域”风险；Depth 的目标是降低某一候选机制的不确定性；Synthesis 用于更新问题图；Experiment 用于把文献主张变成自己的可证伪证据。切换不应按固定轮数，而应由边际信息增益驱动：
\begin{equation}
 m_k^*=\arg\max_m\frac{\E[\Delta I_{decision}\mid m,H_k]}{\E[\Cost(m)\mid H_k]}.
\end{equation}
其中 $I_{decision}$ 不是“知道了多少论文”，而是对后续架构选择、实验设计或动作价值判断真正有影响的信息。由此可得到一个实践停止规则：若同一深挖分支连续新增文献只重复已有机制、不能改变候选集/预测/实验设计，则返回 Breadth；若 Breadth 连续产生的新领域无法改变当前因果图，则进入 Depth 或 Experiment。

\section{RP1 深挖：Temporal-Semantic Retrieval 不是一个检索任务，而是五类任务}

\subsection{从 feature location 到 semantic history slicing}

传统 feature/concept location 主要问“功能现在在哪里”；semantic history slicing 则直接处理“高层功能由哪些历史 change sets 实现”。CSlicer/Definer 系列明确指出，线性 commit history 缺少高层 semantic structure，并以测试/行为作为 slicing criterion，寻找保持目标功能的最小历史子集\parencite{li2018history,li2022semantic,zhu2020genslice}。这与我们的历史检索问题高度相关：研究者真正想找的常常不是某个字符串第一次出现，而是\emph{某个当前语义性质是怎样被一组跨 commit 的变化建立、破坏、修复并最终冻结的}。

因此 repository retrieval 至少应区分：
\begin{align}
R_{now}&: q\to \text{current symbols/files/dependencies},\\
R_{when}&: q\to \text{change points/commits},\\
R_{why}&: q\to \text{decision/evidence episodes},\\
R_{cause}&: q\to \text{minimal behavior-preserving history slice},\\
R_{fail}&: q\to \text{prior failed attempts/counterexamples}.
\end{align}
当前语义 code search、Git search、semantic history slicing 各覆盖其中一部分，尚没有理由相信一个统一 embedding index 能同时最优解决五类任务。工程上应使用 query router，而不是 one-index-to-rule-them-all。

\subsection{检索对象应从 document/snippet 升级为 evidence subgraph}

对研究系统，单个 snippet 往往不是充分证据。一个答案可能需要同时包含：当前代码、引入它的 commit、当时的实验、被否决的替代方案、后续 regression fix。因此更自然的输出是
\begin{equation}
 G_q=(V_q,E_q)\subseteq G_{history},
\end{equation}
即满足 query 的最小充分 evidence subgraph。可定义：
\begin{equation}
G_q^*=\arg\min_{G\subseteq G_H}\Cost(G)
\quad\text{s.t.}\quad
\Pr(AnswerCorrect\mid G,q)\ge 1-\epsilon.
\end{equation}
这把 retrieval 与 provenance 直接接起来：semantic search 负责提出候选节点，结构/时间/因果边负责扩展证据，最终由任务相关 verifier 或 answer consistency 判断是否充分。

\subsection{一个可立即实施的 RP1 benchmark}

不要先发明新 retriever。先从真实 repo 历史构造 30--100 个 archaeology queries，每个 query 冻结 gold evidence set 或至少 expert-validated evidence set。按五类任务分层。比较：
\begin{enumerate}
\item Git/rg baseline；
\item lexical/BM25；
\item semantic code search；
\item symbol/dependency graph navigation；
\item history-aware search；
\item hybrid routed retrieval；
\item agentic multi-step retrieval。
\end{enumerate}
指标不只 Recall@k，而应包含 evidence completeness、time/cost-to-sufficient-evidence、false causal attribution、以及下游 decision accuracy。若 hybrid router 未显著超过简单工具，则不应继续构造复杂历史图。

\section{RP2 深挖：Event Abstraction 是历史系统真正的瓶颈之一}

Process Mining 的 event abstraction 文献明确把问题定义为从 fine-grained event sequence 映射到更粗粒度、适合分析的 event sequence，并指出 mixed granularity 会让发现的过程模型过度复杂、遮蔽真实业务逻辑\parencite{vanzelst2021abstraction,choi2024abstraction}. 这与 agent trace 几乎同构：低层 trace 是 token/model/tool/file/commit 事件，高层研究动作却是 TestHypothesis、AcquireEvidence、RejectMechanism、BranchDirection 等。

设低层事件流
\begin{equation}
E^L=(e_1,\ldots,e_n),
\end{equation}
高层研究事件流
\begin{equation}
E^H=(\bar e_1,\ldots,\bar e_m),\qquad m\ll n,
\end{equation}
则 abstraction 不只是分类：它同时要估计 boundary、grouping、type、参与 objects 与 provenance：
\begin{equation}
\phi:E^L\to (E^H,\Gamma),
\end{equation}
其中 $\Gamma$ 保存每个高层事件到低层证据的映射。

关键约束应为：
\begin{align}
\text{Coverage: }&\quad \forall \bar e,\;Sources(\bar e)\neq\varnothing,\\
\text{Traceability: }&\quad \bar e\to e_{i:j}\text{ 可逆定位},\\
\text{Object preservation: }&\quad Objects(\bar e)\supseteq Objects_{decision-relevant},\\
\text{Temporal consistency: }&\quad order(E^H)\text{ 不制造伪因果},\\
\text{Decision sufficiency: }&\quad \mathcal L_{DR}(E^H)\text{ 足够小}.
\end{align}

\subsection{为什么 object-centric 是必要而非装饰}

2026 年 object-centric process management manifesto 将 object-centricity 描述为对多个共同演化数据对象及其行为相互影响的建模；传统单 case trace 无法自然表达这种结构\parencite{seidel2026ocpm}. 对我们的研究过程，一个低层事件可以同时属于 AgentRun、Candidate、Hypothesis、Experiment、Artifact、GitBranch、Task。强行选择唯一 case ID 会造成 convergence/divergence distortion。因此 ledger 应原生允许 event--object 多对多关系，而不是事后从字符串推断。

\subsection{RP2 最小实验}

从已有真实 campaign 中人工标注一小批高层 research episodes，构成 event-abstraction gold set。比较 rule-based、sequence labeling、change-point detection、LLM segmentation、hybrid。评价：boundary F1、type F1、object-link F1、provenance completeness，以及最重要的 downstream retrieval/forecast delta。若一个 abstraction 方法 summary 很漂亮但让未来检索/预测变差，应判为失败。

\section{RP3 深挖：Memory 应被评价为 state estimator，而不是文本压缩器}

近期 agent memory 已明显从 flat semantic store 转向 event/state/provenance。CompassMem 把 memory 组织为 Event Graph 并沿显式逻辑关系导航\parencite{hu2026compass}; MAGE 明确批评仅按 semantic similarity 组织历史会破坏 execution-state dependencies，转而维护 hierarchical execution-state tree\parencite{chen2026mage}; TrajWiki 保存 immutable episodic snapshots 和 ADD/REVISE/DEPRECATE 演化轨迹\parencite{sun2026trajwiki}; Agent Zero Memory 同时维护 event timeline、entity-event graph 与 citation-locked documentary memory\parencite{wu2026agentzero}. 这些工作共同说明：memory 的核心已经不是“存多少”，而是能否在未来动作前重建足够正确的 state。

因此建议把 memory benchmark 从 QA accuracy 扩展为：
\begin{equation}
\text{StateError},\quad
\text{DependencyRecall},\quad
\text{ContradictionDetection},\quad
\text{StaleMemoryRate},\quad
\text{DecisionRegret}.
\end{equation}
特别是 research memory 必须区分：事实、belief、hypothesis、decision 与 superseded state；否则旧但语义相似的结论会持续污染当前判断。近期 provenance-aware cascading invalidation 的工作也正面处理“新证据到来后依赖旧 memory 的下游记录如何失效，同时保留历史证据”的问题\parencite{wang2026pacmi}。

\section{RP4 深挖：Forecast 的第一目标不是预测下一句话，而是动作条件下的研究前景}

Predictive Process Monitoring 已形成 next activity、remaining time、outcome prediction 等成熟任务。2026 年受控比较显示，在 next-activity prediction 上传统 sequence model 仍稳定优于 LLM，而 tabular foundation model 在时间类任务上有竞争力，LLM 通常成本更高且并不占优\parencite{fertig2026ppm}. 这对我们是强烈警告：不能把 ``LLM world model'' 当默认答案。

研究 forecast 应至少拆成四个 heads：
\begin{align}
\hat p_{imp}(a)&=\Pr(\text{meaningful improvement within }\Delta B\mid S_t,a),\\
\hat g(a)&=\E[\Delta J\mid S_t,a],\\
\hat c(a)&=\E[\Cost\mid S_t,a],\\
\hat i(a)&=\E[\Delta I_{decision}\mid S_t,a].
\end{align}
最终控制器使用的是 action-conditioned distribution，而不是单点 ``next action''：
\begin{equation}
\hat Q(S_t,a)=F(\hat p_{imp},\hat g,\hat c,\hat i,\text{risk}).
\end{equation}

第一代 benchmark 应先用历史 campaign 做离线预测，严格按 campaign/time split，防止同一研究 lineage 泄漏。模型顺序应从 empirical hazard / survival、logistic/GBDT、sequence model、tabular FM，再到 LLM/hybrid；只有复杂模型在 calibration 与 decision utility 上显著胜出才升级。

\section{RP5 深挖：Timing 的理论对象是 costly information + metareasoning，不是 plateau detector}

Monahan 对 partially observable Markov process 中 costly information 与 optimal stopping 的经典分析已经给出一个重要反例：最优购买信息区域不必凸，信息价值关于当前 knowledge state 可以非单调\parencite{monahan1980costly}. 因此“越 plateau 越应检索”只能作为 baseline，不应被写成结构性定理。

更一般地，把研究动作分成内部计算、外部信息获取、实验、人类咨询与停止：
\begin{equation}
\A_M=\{I,E,X,H,Stop\}.
\end{equation}
每个动作都会改变 belief state 和剩余预算。真正目标为
\begin{equation}
\pi^*=\arg\max_\pi\E[U(Z_T)-\sum_t C(a_t)].
\end{equation}
这与 rational metareasoning 的核心一致：计算本身也是有成本的动作，只有其预期决策收益超过成本才值得执行\parencite{russell1991metareasoning}。

\subsection{Human-in-loop 应与 retrieval timing 放在同一 action family}

Sequential Learning-to-Defer 已明确指出逐实例、myopic deferral 忽略长期后果，并用 model-based RL 比较“现在 defer”与“以后 defer”的长期价值\parencite{joshi2021sequentialdefer}. 这与 SearchWorld 的时机问题结构相同。因此 AskHuman 不应是异常处理，而应是昂贵信息/控制动作：
\begin{equation}
VOH_t=\E[U\mid AskHuman,S_t]-\E[U\mid autonomous,S_t]-C_H.
\end{equation}
未来甚至可统一研究 SearchWorld、SpawnSubagent、AskHuman、RunExpensiveExperiment 作为不同的信息 acquisition channels。

\section{RP6 深挖：Taste 不应被神秘化，也不应被过早蒸馏}

Naturalistic Decision Making 把专家直觉解释为长期经验形成的大量 pattern 与 tacit knowledge，而不是通用的简单 heuristic\parencite{klein2015ndm,patterson2010tacit}. Structured Expert Judgment 又提供另一条重要约束：专家意见在数据不足时不可避免，但专家本身必须被 calibration，而不能因为 ``expert'' 身份自动获得权威；Cooke classical model 使用 seed/calibration questions 对专家概率判断进行经验校准\parencite{colson2018expert}。

因此 human taste ledger 至少应记录：
\begin{equation}
h_i=(S_i,A_i^{agent},A_i^{human},R_i,C_i,Y_{future}),
\end{equation}
其中 $R_i$ 是理由，$C_i$ 是人类置信度。长期可估计：
\begin{align}
Calibration_H&=P(Y\mid C_H),\\
OverrideValue&=U(A_H)-U(A_A),\\
DomainOfExpertise&=\{S: \E[OverrideValue\mid S]>0\}.
\end{align}
只有反复、可校准、跨实例稳定的 judgment 才有资格晋升为 SOP；只有具有清晰 mechanical predicate 且 counterexample risk 足够低的 SOP 才晋升为 invariant。否则“蒸馏 taste”会把局部经验误固化成系统性偏差。

\section{一个新的统一对象：Epistemic Control Graph}

第三轮后，单纯的 event log 已不足以表达我们真正要控制的对象。建议逻辑上定义一个可派生、非权威的 Epistemic Control Graph（ECG）：
\begin{equation}
G_t=(V_t,E_t),
\end{equation}
节点类型至少包括
\begin{equation}
V=\{Evidence,Claim,Hypothesis,Decision,Artifact,Experiment,AgentRun,StateView\},
\end{equation}
边类型包括
\begin{equation}
E=\{supports,contradicts,derivedFrom,implementedBy,testedBy,supersedes,spawnedBy,usedBy\}.
\end{equation}
注意：ECG 不是新的 source of truth。权威仍是 append-only raw ledger；ECG 是可重建 projection。其价值在于把 RP1--RP6 接到同一接口：
\begin{align}
Retrieval &: q,G_t\to G_q,\\
State &: G_t\to S_t,\\
Forecast &: S_t,a\to \hat Q,\\
Control &: \hat Q,B_t\to a_t,\\
Human &: (S_t,\hat Q)\to override/advice,\\
Learning &: outcome\to update(views,models,policies).
\end{align}

\section{第三轮提出的三个可证伪核心假设}

为了避免研究继续扩散成 infrastructure collection，建议把未来原创工作压缩到三个承重假设。

\paragraph{H1: Provenance-aware structured retrieval.}
在相同 evidence corpus 与预算下，temporal/structural/provenance-aware retrieval 相比 flat lexical/semantic retrieval，能显著降低达到 sufficient evidence 的成本并减少错误因果归因。

\paragraph{H2: Decision-sufficient abstraction.}
存在远小于完整历史的 state representation $S_t$，满足
\begin{equation}
\mathcal L_{DR}=I(H_t;Q(A))-I(S_t;Q(A))\le\epsilon,
\end{equation}
并在 downstream forecast/control 上不劣于直接使用更长 raw trajectory。

\paragraph{H3: Adaptive epistemic control.}
在同一 task/model/tool/evaluator/hard budget 下，使用 provenance-grounded state 与 calibrated forecast 的 adaptive policy，在 verified progress--cost frontier 上优于 fixed schedules/always-retrieve/never-retrieve/always-defer 等策略。

这三个假设依次对应：能否找到历史；能否把历史压成决策状态；能否用该状态做更好的时机控制。若 H1 失败，应先停止复杂 memory/state；若 H2 失败，应保留 richer state 而非训练 timing policy；若 H3 失败，则说明 state/forecast 即便准确，也未产生足够控制价值。

\section{实验设计：从现在即可开始积累的四个数据集}

\subsection{D1 Repository Archaeology Set}
从真实长期仓库中抽取 ``where/when/why/cause/fail'' 查询；保存 expert evidence graph、相关 commit/diff/test/decision。它服务 RP1。

\subsection{D2 Research Episode Abstraction Set}
对真实 agent/tool/git trace 标注高层 research episode boundary/type/object links。它服务 RP2，并可评估 event abstraction 对 RP1/RP4 的 downstream 影响。

\subsection{D3 Campaign Forecast Set}
每个 campaign 在多个 checkpoint 冻结 $S_t$，记录后续固定 budget window 内各类动作的实际 outcome/cost。初期可用 observational data，随后需要 randomized interventions 才能可靠估计 action-conditioned value。

\subsection{D4 Human Judgment Calibration Set}
所有 human override/defer/continue/stop/retrieve 建议均记录预测、置信度、理由与后续 outcome；不要只记录最终修改。这是未来 taste calibration 与 judgment crystallization 的唯一低成本原始材料。

\section{因果警告：历史数据不足以直接学习最优研究 policy}

一个严重风险是从历史 trajectories 直接训练 ``best next action''。历史 action 是旧 agent/human policy 的产物：
\begin{equation}
a_t\sim\pi_{old}(a\mid S_t),
\end{equation}
因此 $S_t$ 同时影响 action 与 outcome，且许多未记录因素也影响二者。观察到 ``retrieval 后经常改善'' 不等于 retrieval 的 causal effect 为正。Prescriptive Process Monitoring 和 Learning-to-Defer 的近期研究同样面对 intervention policy 与 hidden confounding 问题\parencite{montreuil2026causaldefer,delias2026policy}。

因此路线应为：
\begin{enumerate}
\item observational prediction：只验证 state/forecast 能否预测；
\item randomized micro-interventions：在低风险动作上随机 timing/route，估计 treatment effect；
\item offline policy evaluation：仅在 overlap 足够、propensity 可估时使用；
\item prospective policy trial：最终比较 policy-level utility。
\end{enumerate}
不能把 predictive accuracy 偷换成 control validity。

\section{工程架构进一步收敛：四个平面，而非 Research OS}

第二轮的 Control/Epistemic 两平面仍可进一步细化为四个相互解耦平面：
\begin{enumerate}
\item \textbf{Execution Plane}：真正执行 agent/tool/model/evaluator/git；
\item \textbf{Evidence Plane}：append-only raw events、artifacts、hash、trace/provenance；
\item \textbf{Epistemic Plane}：可重建的 semantic/event/state/retrieval views；
\item \textbf{Control Plane}：IaC/policy/budget/permissions/evaluator identity 与 metareasoning policy。
\end{enumerate}
Human judgment 横跨 Epistemic 与 Control，但每次介入都必须回写 Evidence Plane。这样任何复杂的 memory graph、forecast model 或 policy 都可以替换，而 raw evidence 与 hard constraints 不受影响。

\section{研究推进 policy：何时继续深挖，何时回到广度}

本轮之后，不建议机械地“先把 RP1 做完再 RP2”。它们存在强耦合：RP1 暴露检索缺失会反向要求 RP2 增加某类 event/object；RP2 的新 abstraction 又会改变 RP1 query。更合理的是小闭环交替：
\begin{equation}
Breadth\to Depth(RP_i)\to Prototype\to Failure\to CrossDomainSearch\to Synthesis\to Depth(RP_j).
\end{equation}

定义研究前沿状态 $F_k$，每次循环只允许因以下原因扩展广度：
\begin{enumerate}
\item 当前模型出现无法解释的 counterexample；
\item 现有候选方法在关键指标上同时失败；
\item 出现接口缺口，且邻接学科可能已有成熟解；
\item 新术语/理论能改变实验设计或可证伪预测。
\end{enumerate}
反之，若新文献只改变名词、不改变 $F_k$ 的候选、预测或实验，则停止该方向横向搜索。深度也有停止条件：当 source-level mechanism、适用边界、失败模式、benchmark、实现成熟度均已足以做 build-vs-reuse 决策时，立即进入 prototype/experiment，不继续“为了完整”读文献。

\section{第三轮结论：总问题的最小研究内核}

经过四轮检索，最小研究内核已可写成：
\begin{equation}
\boxed{
\begin{aligned}
&\textbf{Evidence}: &&\text{发生了什么，证据在哪里？}\\
&\textbf{Abstraction}: &&\text{哪些低层事件构成高层研究状态？}\\
&\textbf{Retrieval}: &&\text{当前决策最小充分证据是什么？}\\
&\textbf{Forecast}: &&\text{采取不同动作后会怎样？}\\
&\textbf{Control}: &&\text{有限预算下现在最值得做什么？}\\
&\textbf{Judgment}: &&\text{机器何时应 defer，人的 taste 如何校准并结晶？}
\end{aligned}}
\end{equation}

因此当前最值得做的不是继续设计一个更大的 agent framework，而是开始积累 D1--D4 四类可复用数据，并用它们逐一攻击 H1--H3。底层尽量采用已有 telemetry/provenance/process-mining/code-search 组件；原创研究只发生在已有组件无法回答的接口上。若未来某个新框架不能明确改善 H1/H2/H3 中至少一个可测量量，则它不应进入主架构。



\part{第四轮：从主问题反推工程与研究内核}

\section{重新冻结主问题：组件不是目标}

前述调研容易出现一种典型偏差：某个子问题（memory、forecasting、swarm、RAG、process mining）文献丰富，遂反过来把系统设计成该子领域的放大版。第四轮明确禁止这种倒置。冻结唯一上位目标
\begin{equation}
\boxed{G_0=\text{使开放世界中的真实算法研究更适合 Agent 自动、可靠、长期地推进。}}
\end{equation}
所谓 ``真实'' 至少意味着：任务和知识边界持续变化；代码库、数据、依赖和模型版本会漂移；研究存在失败、回滚、分叉、合并和人类介入；外部网络和 Git 历史是可用资源；预算是有限的；系统必须在数小时到数周的时间尺度上可恢复，而非仅在一次 prompt 内完成。

因此，本文不把 ``构造 Research OS''、``构造统一知识图谱''、``训练统一 world model'' 或 ``最大化 agent autonomy'' 作为目标。任何组件 $c$ 只有在能回答
\begin{equation}
\Delta G_0(c)>0
\end{equation}
且其收益不能由更成熟、更廉价的 shelf solution 获得时，才进入原创实现范围。

\section{三重困难的动力学：epistemic debt 而非单点失败}

令 $H_t$ 为截至时刻 $t$ 的完整研究历史，$S_t^*$ 为决策所需但不可直接观测的真实研究状态，$\hat S_t$ 为系统重建状态，$a_t$ 为研究动作。理想闭环为
\begin{equation}
H_t\xrightarrow{\phi}\hat S_t\xrightarrow{\pi}a_t\longrightarrow H_{t+1}.
\end{equation}
定义三类误差：检索误差 $\epsilon_R$、状态/控制误差 $\epsilon_S$、历史组织误差 $\epsilon_H$。它们不是独立加性噪声，而存在正反馈：
\begin{equation}
\boxed{
\epsilon_H\uparrow\Rightarrow
\epsilon_R\uparrow\Rightarrow
\epsilon_S\uparrow\Rightarrow
P(a_t\text{ is wasteful})\uparrow\Rightarrow
\epsilon_H\uparrow.}
\end{equation}
这比 ``context 太长'' 更一般。我们称累积的、使未来系统越来越难恢复证据、重建状态和作出正确研究决策的负担为 \emph{epistemic debt}。2026 年关于 long-horizon search 的 context-rot 工作已经表明，仅增加上下文会导致放弃、提前回答和性能下降，并系统比较了多种 context management 策略\cite{xia2026contextrot}；ARC 则把 context management 明确改写为主动的动态内部状态维护\cite{yao2026arc}。这说明 ``把历史都留在 context'' 既不是历史处理，也不是可靠 memory。

\section{由三重困难得到四个系统能力，而非九个孤立模块}

从 $G_0$ 出发，最小能力可压缩为四问：
\begin{enumerate}
\item \textbf{Recoverability}：过去与当前决策有关的证据，能否以低成本、低误归因恢复？
\item \textbf{Situational Awareness}：能否从历史与当前观测形成足够好的研究状态？
\item \textbf{Resource Rationality}：下一单位 token、wall-clock、实验、检索、subagent 或 human attention 应投向哪里？
\item \textbf{Organizational Learning}：过去的经验是否降低未来同类问题的认知与执行成本？
\end{enumerate}
四者构成
\begin{equation}
\boxed{Recover\to Understand\to Allocate\to Learn\to Recover'.}
\end{equation}
Retrieval、event abstraction、context assembly、forecasting、swarm scheduling、SOP、human deferral 均只是实现这四项能力的手段。

\section{Evidence-first：研究过程必须产生可复用证据}

最基础的工程不变量不是 ``所有内容进入 memory''，而是重要研究动作必须留下可寻址证据：
\begin{equation}
\boxed{Action\to Artifact+Evidence+Provenance.}
\end{equation}
例如 subagent 的最终产物不应只有自然语言结论，而至少包含 claim、supporting evidence、artifacts、scope、confidence/provisional status 与 provenance；实验不应只有 stdout，而应绑定 hypothesis、code/config/data、evaluator identity、outcome 与 interpretation；决策应绑定 alternatives、evidence、forecast-at-decision-time、choice 与 reason。

科学工作流领域已经有可直接复用的 provenance 基础。Workflow Run RO-Crate 可以把 workflow run 的 inputs、outputs、code、step executions 和 intermediate outputs 打包，并与 W3C PROV 对齐\cite{soiland2024wfrun}。因此我们的新工作不应重造 generic provenance，而是补上 ``research-semantic event'' 与已有 provenance 的映射。

\section{历史处理：从 telemetry 到 typed temporal research state}

\subsection{冷热分层与晋升边界}
真实长程 agent 会产生大量低层 telemetry。OpenTelemetry 的 GenAI semantic conventions 已经能标准化 model、token、tool-call、latency 等执行观测\cite{otel2026genai,opentelemetry2026inside}。这类数据应尽量直接复用，而不进入原创 ontology。

建议逻辑分层：
\begin{equation}
\boxed{
L_0\ \text{raw telemetry}
\to L_1\ \text{execution trace}
\to L_2\ \text{research events}
\to L_3\ \text{derived state}
\to L_4\ \text{working context}.}
\end{equation}
通常 $|L_0|\gg|L_1|\gg|L_2|\gg|L_4|$。核心研究问题是晋升映射
\begin{equation}
\alpha:\mathcal E_{low}^{*}\to\mathcal E_{research},
\end{equation}
而不是如何采集更多 log。Process Mining 的 event abstraction 文献已经系统研究了 fine-grained events 到 coarse-grained events 的映射，并指出 mixed granularity 会产生过度复杂、掩盖真实业务逻辑的过程模型\cite{vanzelst2021abstraction}。我们的差异在于目标不是 business activity，而是 hypothesis testing、evidence acquisition、branch rejection、artifact validation 等研究语义。

\subsection{Object-centric 而非单 trajectory}
单一 case/trajectory 模型不适合 swarm 与研究 DAG。一个事件可能同时关联 AgentRun、Candidate、Hypothesis、Experiment、Artifact、GitBranch 与 Task。OCEL 2.0 原生支持 event--object、object--object 关系、qualifier 以及随时间变化的 object attributes\cite{berti2024ocel}；2026 年 object-centric process management manifesto 进一步把共同演化对象视为 process modeling/mining 的核心\cite{seidel2026ocpm}。故第一版数据模型应优先借用 object-centric 思想，而非强迫所有事件归属唯一 session。

\subsection{Typed temporal knowledge 与 invalidation}
历史中的对象生命周期不同。至少区分：
\begin{align}
Fact &: \text{immutable observation},\\
Interpretation &: ACTIVE\to SUPERSEDED/DISPUTED,\\
Forecast &: PENDING\to RESOLVED,\\
Decision &: ACTIVE\to REVERSED,\\
Failure &: OPEN\to EXPLAINED\to PREVENTED.
\end{align}
因此 knowledge item 更合理地表示为
\begin{equation}
k_i=(claim,source,t_{valid},scope,confidence,status,provenance).
\end{equation}
否则 retrieval 可能稳定地返回已经失效但语义高度相似的旧结论。

\subsection{Branch 与 Failure 必须是一等对象}
真实研究是 DAG 而非单链：
\begin{equation}
ResearchGraph=(V_{hyp}\cup V_{branch}\cup V_{artifact}\cup V_{exp},E).
\end{equation}
Branch 至少保存 goal、parent、state、artifacts、evidence、cost 与 status；Failure 至少保存 context、attempt、symptom、mechanism/root-cause hypothesis、evidence、scope、resolution。若失败只存在于聊天文本，系统无法可靠避免重复失败，organizational learning 也无从发生。

\section{检索：从 ``找 chunk'' 到 ``恢复最小充分证据子图''}

\subsection{检索任务必须拆分}
软件工程中的 feature location 已经形成庞大文献；系统综述覆盖了 89 篇工作，并表明文本、结构、动态与历史信号是不同方法族\cite{dit2013feature}。Semantic History Slicing 更进一步：目标不是关键词命中，而是找出实现目标功能的、语义相关且可执行的历史 changeset 子集\cite{li2022semantic}，CSlicer/DEFINER 提供了原型实现\cite{cslicerrepo}。

对研究系统，至少区分：
\begin{equation}
Q\in\{where,when,why,cause,fail,expected\}.
\end{equation}
其中 $why$ 需要跨 Decision/Experiment/AgentTrace，$cause$ 需要结构与历史依赖，$expected$ 需要恢复决策时的 forecast，而不是当前事后解释。故不存在单一 embedding retriever 可自然解决全部问题。

\subsection{真正的检索输出}
定义历史图 $G_H$。对于查询 $q$，理想输出不是 top-$k$ chunks，而是最小充分证据子图
\begin{equation}
G_q^*=\arg\min_{G'\subseteq G_H}\Cost(G')
\quad\text{s.t.}\quad
P(\text{correct decision}\mid G',q)\ge 1-\epsilon.
\end{equation}
这是一个目标定义，不要求第一版直接求最优。工程上可采用 exact-ID $\to$ structural $\to$ temporal $\to$ lexical/semantic $\to$ external 的 routed retrieval，再由 context assembler 做裁剪。

\subsection{历史可寻址性比向量召回更基础}
重要对象应有跨时间 stable identity：
\begin{equation}
M_{17}^{(0)}\to M_{17}^{(1)}\to M_{17}^{(2)},
\end{equation}
即使文件路径、函数名和实现语言改变，Mechanism $M_{17}$ 的演化仍可追踪。工程第一版只需要 stable IDs、typed relations、timestamps、hashes 与 provenance；并不需要先部署庞大知识图谱。

\section{Context Assembly：Storage、Retrieval 与 Working Context 必须分离}

系统拥有证据不等于模型此刻应该看到全部证据。定义角色 $r$、任务 $g_t$、状态 $S_t$，working context 为
\begin{equation}
K_t=C(g_t,S_t,E_t,R_t,B_t,r),
\end{equation}
其中 $E_t$ 为检索证据，$R_t$ 为有效规则。优化目标不是 context 最大，而是
\begin{equation}
\min |K_t|\quad\text{s.t.}\quad K_t\text{ 对当前决策充分}.
\end{equation}
MAGE 将 memory 重新表述为 execution-state management，用层次状态树维护 active root-to-current path，并通过 Grow/Compress/Maintain/Revise 隔离错误分支\cite{chen2026mage}；CompassMem 则用 event graph 和显式逻辑关系支持结构化导航\cite{hu2026compass}。这些结果共同支持 ``history $\neq$ memory $\neq$ context''。

对 swarm，context allocation 本身还是实验设计变量：若希望两个 subagents 独立探索，给二者完全相同的历史可能降低 diversity。故 scheduler 实际决定
\begin{equation}
(Who,WhatTask,WhatContext,WhatBudget),
\end{equation}
而非仅仅 $Spawn(N)$。

\section{状态、预测与时机：统一成有限资源下的 epistemic control}

\subsection{状态估计}
研究状态 $S_t$ 不应只是摘要，而应至少包含 goal、active hypotheses、accepted/rejected evidence、open failures、active branches、pending experiments、uncertainties、budgets 与 active policies。目标是寻找压缩表示
\begin{equation}
S_t=\phi(H_t),\qquad |S_t|\ll |H_t|,
\end{equation}
同时尽量保留未来决策相关信息。可用概念量
\begin{equation}
\mathcal L_{DR}=I(H_t;Q(A))-I(S_t;Q(A))
\end{equation}
描述 decision-relevant information loss；实际实验中用 downstream regret、duplicate work、invalidated-knowledge reuse 等可测代理量替代直接互信息估计。

\subsection{预测必须多时间尺度，但服务于资源控制}
一步预测对下一 action 有直接价值；多步/结构预测对 option、branch、subagent 和项目方向有粗粒度价值。可表示为
\begin{equation}
\mathcal F=\{F_{step},F_{segment},F_{branch},F_{struct},F_{resource},F_{terminal}\}.
\end{equation}
长 horizon 不必精确预测 microstate，而应降低表示分辨率：
\begin{equation}
h\uparrow\Rightarrow resolution(\phi_h)\downarrow.
\end{equation}
最终评价预测器不是只看 predictive loss，而看
\begin{equation}
V_{pred}=U(\pi_{\hat F})-U(\pi_{no\ forecast}).
\end{equation}
即它是否真的改善研究决策。

\subsection{时机的主问题是 ``下一单位研究资源投哪里''}
动作空间可以统一为
\begin{equation}
\begin{gathered}
\A_M=\{InternalCompute,Code,Experiment,RetrieveLocal,RetrieveWorld,\\
SpawnAgent,Continue,Prune,Merge,AskHuman,Stop\}.
\end{gathered}
\end{equation}
控制目标为
\begin{equation}
a_t^*=\arg\max_a \E[\Delta V_{research}\mid S_t,a,B_t].
\end{equation}
Prescriptive Process Monitoring 是比单纯 Active RAG 更直接的邻接领域：它研究 ongoing case 上何时 intervention，并把预测、不确定性、因果效应与有限资源共同纳入 allocation\cite{shoush2022resource,shoush2022when,delias2026policy}。这提醒我们：预测 ``哪条 branch 会失败'' 不等于知道 ``现在干预该 branch 是否有因果收益''。

经典 costly-information POMDP 甚至表明最优购买信息区域不必是凸的，信息价值随当前 knowledge state 可非单调\cite{monahan1980costly}。因此 ``plateau 越严重越应该检索'' 只能是 baseline，而非一般理论。

\section{IaC、SOP 与 Taste：三种知识成熟度，而非三种文档格式}

\subsection{Executable epistemic constraints}
凡可机械验证的约束，应尽量下沉到 code/policy：artifact hash、evaluator identity、budget、branch parent、deprecated claim status、commit/run binding 等。原则为
\begin{equation}
CanBeMechanicallyVerified\Rightarrow MustNotRemainNaturalLanguageOnly.
\end{equation}
这里 IaC 应广义理解为 executable epistemic constraints，而不是特定基础设施工具。

\subsection{SOP 是启发式 policy}
SOP 本质是
\begin{equation}
\pi_{SOP}(S)\to RecommendedAction,
\end{equation}
Markdown 只是最低成本表示。高频、稳定、可判别的 SOP 可逐步转为 structured rule / policy-as-code；低频且依赖语境的规则仍保留自然语言。

\subsection{Taste 是 open-set residual controller}
设所有研究决策机会为 $D_0$，可机械解决者为 $D_{hard}$，稳定 procedure 可覆盖者为 $D_{proc}$，则
\begin{equation}
D_{taste}=D_0-D_{hard}-D_{proc}.
\end{equation}
开放世界持续产生 out-of-distribution states，因此无需假设 $D_{taste}\to\varnothing$。Naturalistic Decision Making 把专家直觉与长期经验形成的 patterns/tacit knowledge 联系起来\cite{klein2015ndm}；Structured Expert Judgment 又提醒我们专家判断必须校准，而不能仅凭身份赋予权威\cite{colson2018expert,colson2017cross}。

因此 HumanOverride 应保存
\begin{equation}
h_i=(S_i,A_i^{agent},A_i^{human},Reason_i,Forecast_i,Confidence_i,Outcome_i),
\end{equation}
以后才能估计 human calibration、override value 与 domain of expertise。

\subsection{知识结晶是双向状态机}
长期学习不是无限增加 MD，而是
\begin{equation}
\boxed{Judgment\rightleftarrows Procedure\rightleftarrows Invariant.}
\end{equation}
Promotion 需要独立重复证据、收益、scope 和反例检查；出现环境漂移或 counterexample 时允许 demotion。若同类 human override 反复发生却从不转化为 procedure/invariant，系统并没有真正组织学习。

\section{工程架构：四平面与双速度证据闭环}

推荐逻辑架构：
\begin{enumerate}
\item \textbf{Execution Plane}：真正执行模型、工具、实验、Git、workflow；
\item \textbf{Evidence Plane}：append-only raw/structured evidence、artifact hash、provenance；
\item \textbf{Epistemic Plane}：可重建的 state、retrieval index、failure views、forecast views；
\item \textbf{Control Plane}：hard constraints、SOP/policy、budget、permissions、scheduler、human deferral。
\end{enumerate}
Human judgment 横跨 Epistemic/Control，但所有介入必须写回 Evidence。重要原则是 ``facts first, views disposable''：换 retriever、graph schema、forecast model 或 agent runtime，不应破坏底层事实。

运行反馈至少有三个时间尺度：
\begin{equation}
\tau_{action}<\tau_{research}<\tau_{governance}.
\end{equation}
Fast loop 做工具/代码/测试；medium loop 做 branch、retrieve、allocate、stop；slow loop 才做 SOP/policy/crystallization。避免每个 fast-loop outcome 都立即修改系统规则，否则 attribution 与稳定性会丢失。

\section{Shelf-first 边界：哪些坚决不重新发明}

现有成熟或快速成熟能力应作为依赖/适配层：
\begin{itemize}
\item telemetry/trace：OpenTelemetry GenAI conventions\cite{opentelemetry2026inside};
\item generic provenance/package：W3C PROV / Workflow Run RO-Crate\cite{soiland2024wfrun};
\item object-centric event representation：OCEL 2.0\cite{berti2024ocel};
\item Git/version history：Git + existing MSR/semantic history tools\cite{li2022semantic};
\item generic durable workflow、SQL/Parquet、full-text/vector indexes：直接采用成熟工程组件。
\end{itemize}
原创 intellectual core 只保留：ResearchEventAbstraction、Temporal-Semantic Retrieval、Decision-Relevant Context Assembly、Research State Estimation、Research Foresight、Epistemic Action Selection、Branch/Swarm Value Allocation、Experience Crystallization、Human Deferral。

\section{最小可实施原型：不是 Research OS}

第一版只需：
\begin{enumerate}
\item stable IDs；
\item typed research events；
\item artifact/provenance links；
\item CurrentState projection；
\item temporal/structural retrieval；
\item HumanOverride recording；
\item FailureEvent recording。
\end{enumerate}
物理实现完全可以是
\begin{equation}
\boxed{SQLite/Parquet+Git+filesystem+existing\ telemetry.}
\end{equation}
建议最小对象集合：
\begin{equation}
\begin{gathered}
\{Task,Goal,Hypothesis,Branch,AgentRun,Experiment,Artifact,Claim,\\
Decision,Forecast,Failure,Rule\}.
\end{gathered}
\end{equation}
所有对象只要求 stable ID；关系采用小集合如 $producedBy,usedBy,tests,supports,\allowbreak contradicts,supersedes,\allowbreak derivedFrom,parentOf,resolvedBy$。不要在 MVP 阶段追求完备 ontology。

\section{评价体系：从组件准确率转向研究闭环效用}

\subsection{Recoverability}
定义查询 $q$ 下恢复足够证据的成本
\begin{equation}
CSE(q)=\Cost(RecoverSufficientEvidence\mid q).
\end{equation}
长期目标不是历史变小，而是
\begin{equation}
\frac{d\E[CSE]}{d|H|}\approx 0.
\end{equation}
同时测 FalseCausalAttribution、StaleEvidenceReuse 与 DecisionContextRecovery。

\subsection{State quality}
从 campaign checkpoint 隐藏后续历史，让新 agent 仅通过系统恢复状态，测：
\begin{equation}
DuplicateWorkRate,\ InvalidatedKnowledgeReuse,\ OpenQuestionRecall,\ VerifiedProgressAfterResume.
\end{equation}
这比 summary ROUGE 或 memory QA 更贴近 $G_0$。

\subsection{Resource rationality}
主指标为固定 hard budget 下
\begin{equation}
V_B(\pi)=\E[J_{best}(B;\pi)-J_0],
\end{equation}
并报告 token、wall-clock、evaluator calls、human minutes、branch duplication 与 failure rate。Forecast 单独报告 calibration/predictive accuracy，但最终必须报告 DecisionValue。

\subsection{Organizational learning}
定义重复 judgment class 的结晶率
\begin{equation}
CR=\frac{\#\text{repeated judgment classes successfully crystallized}}{\#\text{repeated judgment classes}},
\end{equation}
同时要求 error rate 不上升。还应测 HumanLoad/VerifiedProgress 随时间是否下降。

\section{Field 与 Lab：发现问题和识别因果必须双轨}

仅实验室控制变量会错过真实 repository drift、network/search noise、tool failure、human interruption、model update 与 context pollution；仅 field success 又无法识别机制。故采用
\begin{equation}
\boxed{Field\to Failure\to Lab\to Mechanism\to Policy\to Field.}
\end{equation}
Field track 负责发现 failure modes 和真实成本分布；Lab track 负责 randomized micro-intervention、ablation 与 causal identification。Prescriptive process monitoring 的综述同样指出 causality、second-order effects 和真实环境验证是关键缺口\cite{fertig2026ppm}。

现有项目本身即可作为 field site。每次重复旧实验、错误恢复历史、stale context、subagent 重复劳动、过早/过晚检索、SOP 未遵守、人类纠正 agent、意外 prior-art discovery，都记录为 Failure/Intervention event。这样 ontology 由真实失败逐渐生长，而不是先验设计。

\section{四个正式研究问题与 kill chain}

\paragraph{RQ1 Recoverability.}
结构化 provenance/temporal/structural retrieval 是否在相同 corpus 与 budget 下显著降低 $CSE$、错误归因和 stale-evidence reuse，相比 flat lexical/semantic retrieval？

\paragraph{RQ2 Decision-sufficient state.}
是否存在 $|S_t|\ll|H_t|$ 的 compact state，使 downstream decision regret、重复工作率和恢复后 verified progress 与 full-history oracle 接近？

\paragraph{RQ3 Foresight and resource rationality.}
多尺度 forecast 是否具有正的 decision value；并且 adaptive resource policy 是否优于 fixed schedule / greedy-myopic policy？注意 predictive accuracy 不足以证明控制价值。

\paragraph{RQ4 Organizational learning.}
记录并结晶 HumanOverride/Failure 是否能在不提高错误率的前提下降低同类 future human load 与 recovery/control cost？

对应 kill chain：
\begin{equation}
\boxed{RQ1\to RQ2\to RQ3,\qquad RQ4\text{ parallel and longitudinal}.}
\end{equation}
若 RQ1 FAIL，复杂历史结构可能不值得；RQ1 PASS 而 RQ2 FAIL，则结构化历史有用但不能安全压缩为 state；RQ1--RQ2 PASS 而 RQ3 FAIL，则 state/forecast 没有足够控制价值。RQ4 需要长期 field data，不能被一次短 benchmark 替代。

\section{正式结论：主问题的当前最小理论}

第四轮后，最稳健的统一表述不是 ``Research State Estimation and Control'' 本身，而是：
\begin{equation}
\boxed{\begin{gathered}
\textbf{Long-horizon research as an epistemic organization}\\
\textbf{under partial observability and bounded resources.}
\end{gathered}}
\end{equation}
其中研究组织由 Human、LLM/Agent、Tool、Repository、Experiment 和 External World 共同组成。系统必须持续解决四件事：
\begin{equation}
\boxed{
\begin{aligned}
Recover &: \text{以低成本恢复与当前决策有关的过去；}\\
Understand &: \text{形成可追溯、可失效、可更新的当前研究状态；}\\
Allocate &: \text{把下一单位有限资源分配到最高预期研究价值处；}\\
Learn &: \text{把重复经验转化为可复用 evidence/procedure/invariant。}
\end{aligned}}
\end{equation}

这也重新解释最初的三层结构：
\begin{equation}
\boxed{Taste\rightleftarrows SOP\rightleftarrows IaC/Invariant.}
\end{equation}
它们不是三个独立存储层，而是同一知识从开放、依赖经验的 judgment 到可描述 procedure，再到可机械执行 invariant 的成熟度轴。Human-in-the-loop 的长期目标也不是 ``永远审批每一步''，而是作为 open-set residual controller，在 novel/high-impact/low-confidence/conflict/irreversible states 中介入，并让每次介入成为可校准、可结晶的数据。

最终工程原则可以压缩为九条：
\begin{enumerate}
\item \textbf{Goal-first}：任何组件必须证明其与 $G_0$ 的关系；
\item \textbf{Evidence-first}：重要动作产生 artifact+evidence+provenance；
\item \textbf{Facts immutable, views disposable}：历史不被 summary 替代；
\item \textbf{Object-centric}：研究是多对象共同演化，而非单 trajectory；
\item \textbf{Storage $\neq$ retrieval $\neq$ context}：三者分别优化；
\item \textbf{Prediction serves allocation}：forecast 以 decision value 而非单纯 accuracy 评价；
\item \textbf{Shelf-first}：telemetry、workflow、database、generic provenance 不重造；
\item \textbf{Field $\leftrightarrow$ Lab}：field 发现失效模式，lab 识别因果；
\item \textbf{Crystallize conservatively}：Judgment/SOP/Invariant 可升可降，且必须保留 scope、provenance 与 counterexample。
\end{enumerate}

因此，当前最合理的下一步不是继续扩张架构，而是让真实研究运行开始产生 RQ1--RQ4 所需的数据，并首先实现一个极薄的 Evidence/State prototype。若该薄层不能在真实 campaign 中降低恢复成本、重复劳动或 human load，就应当缩减甚至放弃，而不是用更大的 ontology、更多 agent 或更复杂 world model 掩盖失败。这是从 $G_0$ 出发最重要的可证伪约束。

\printbibliography
\end{document}
